% !TeX root = ../../Book5.tex
% 纲要标题：## 第13章　Killing 型与半单 Lie 代数
% 纲要位置：工作记录/计划纲要.md，第 4482 行。

\chapter{Killing 型与半单 Lie 代数}
\label{b5:ch:13}
\BookChapterNumber{b5:ch:13}

\begin{chapterabstract}
	本章证明 Cartan 的可解性与半单性判据，由 Killing 型的非退化性分解单理想，并在经典矩阵 Lie 代数中检验这些结构。
\end{chapterabstract}

\begin{learninggoals}
	\begin{itemize}
		\item 能从迹的循环性证明不变性，区分 Killing 型与其他表示迹型。
		\item 用谱插值将 Cartan 判据中的迹条件变为幂零性。
		\item 能用正交补分解半单 Lie 代数的理想，并解释具体单理想的唯一性。
		\item 计算经典矩阵 Lie 代数的维数和中心，证明其单性范围并解释低维例外。
		\item 区分 Lie 代数同构与相应 Lie 群同构，不把有限中心或覆盖问题混入线性层面的结论。
	\end{itemize}
	预备知识是\chapref{b5:ch:11}的伴随表示和理想、\chapref{b5:ch:12}的 Engel 与 Lie 定理，以及线性代数的广义特征空间、双线性型和外幂。
\end{learninggoals}

在 \(\mathfrak{sl}_2(\mathbb C)\) 中，取 \(e=E_{12}\)、\(f=E_{21}\)、\(h=E_{11}-E_{22}\)。伴随作用的迹配对 \(\kappa(x,y)=\operatorname{tr}(\operatorname{ad}x\operatorname{ad}y)\) 给出 \(\kappa(h,h)=8\)、\(\kappa(e,f)=4\)，所以 Killing 型非退化。若只取上三角子代数 \(\langle h,e\rangle\)，其 Killing 型却在 \(e\) 方向退化。迹型如何检出可解理想，非退化性又为什么足以迫使整个 Lie 代数半单？本章从这些矩阵计算进入 Cartan 判据，随后以正交补构造半单代数的单理想分解。有关经典证明可对照\cite[Sections 4.3, 5.1--5.2, pp.~19--23]{Humphreys1972LieAlgebras}。

% !TeX root = ../../Book5.tex
% 纲要标题：### 13.1 不变双线性型
% 纲要位置：工作记录/计划纲要.md，第 4484 行。

\section{不变双线性型}
\label{b5:sec:1301}

% 纲要范围：
%
% * Killing 型
% * 不变性及根空间正交的动机
% * 单 Lie 代数上的形式
% * 比较自然迹型与 Killing 型，通过伴随对偶解释相反权配对，并说明不变型唯一性的基域条件
%

上一章用上三角矩阵识别可解性。本章通过伴随算子的乘积取迹，构造一个双线性型。Cartan 判据将证明，在有限维复 Lie 代数中，这个型非退化恰好等价于 Lie 代数半单。

\subsection{Killing 型}
\label{b5:subsec:1301:01}

\begin{definition}\label{b5:def:killing-form}
	设 \(\mathfrak g\) 是域 \(k\) 上的有限维 Lie 代数。映射
	\[
		\kappa_{\mathfrak g}(x,y) =\operatorname{tr}_{\mathfrak g} \left(\operatorname{ad}x\circ\operatorname{ad}y\right), \qquad x,y\in\mathfrak g,
	\]
	称为 \(\mathfrak g\) 的\term[Killing-xing]{Killing 型}[Killing form]。更一般地，若 \(V\) 是有限维 \(k\) 向量空间，表示 \(\rho:\mathfrak g\rightarrow\mathfrak{gl}(V)\) 给出\term[ji-xing]{迹型}[trace form] \(\beta_\rho(x,y)=\operatorname{tr}_V\bigl(\rho(x)\rho(y)\bigr)\)。
	\symidx{\(\kappa_{\mathfrak g}\)：Lie 代数的 Killing 型}
\end{definition}
\symidx{\(\beta_\rho\)：有限维 Lie 代数表示的迹型}

\begin{example}\label{b5:exm:killing-sln}
	对 \(X,Y\in\mathfrak{sl}_n(\mathbb C)\)，
	\begin{equation}\label{b5:eq:killing-sln}
		\kappa_{\mathfrak{sl}_n}(X,Y)=2n\operatorname{tr}(XY).
	\end{equation}
	为核验系数，在 \(M_n(\mathbb C)\) 上记 \(L_X(T)=XT\)、\(R_Y(T)=TY\)。矩阵单位基给出
	\[
		\operatorname{tr}(L_A)=\operatorname{tr}(R_A)=n\operatorname{tr}A, \qquad \operatorname{tr}(L_XR_Y)=\operatorname{tr}X\,\operatorname{tr}Y.
	\]
	最后一式中 \(E_{ij}\) 的自身系数为 \(X_{ii}Y_{jj}\)，求和即可。展开 \((L_X-R_X)(L_Y-R_Y)\)，得
	\[
		\kappa_{\mathfrak{gl}_n}(X,Y) =2n\operatorname{tr}(XY)-2\operatorname{tr}X\,\operatorname{tr}Y.
	\]
	又因 \(\mathfrak{gl}_n=\mathfrak{sl}_n\oplus\mathbb CI\)，伴随作用在标量直线上为零，便得到\eqref{b5:eq:killing-sln}。特别在 \(n=2\) 时，对标准的 \(e=E_{12}\)、\(f=E_{21}\)、\(h=E_{11}-E_{22}\)，有
	\[
		\kappa(h,h)=8,\qquad \kappa(e,f)=4,\qquad \kappa(e,e)=\kappa(f,f)=\kappa(h,e)=\kappa(h,f)=0.
	\]
	因此它在基 \(h,e,f\) 下的矩阵为
	\[
		\begin{pmatrix}8&0&0\\0&0&4\\0&4&0\end{pmatrix},
	\]
	行列式为 \(-128\ne0\)，故这个 Killing 型非退化。
\end{example}

这个计算也区分了 Killing 型与一般表示的迹型。对 \(n\ge2\)，\(\mathfrak{gl}_n(\mathbb C)\) 的自然表示给出 \(\beta_\rho(X,Y)=\operatorname{tr}(XY)\)。若它对固定的 \(X\) 与所有 \(Y\) 的配对都为零，逐一取 \(Y=E_{ji}\)，便得 \(X_{ij}=0\)，所以这个迹型非退化。但 \(\mathbb CI\) 是非零交换理想，\(\mathfrak{gl}_n\) 因而不半单；又因 \(\operatorname{ad}I=0\)，有 \(\kappa_{\mathfrak{gl}_n}(I,Y)=0\) 对所有 \(Y\) 成立，Killing 型退化。一般表示可以让中心元素非零地作用，伴随表示却将中心零化。Cartan 半单性判据所用的是这个内部伴随作用的迹型，不能只用某个表示的迹型非退化来判定半单性。

\subsection{不变性与根空间正交}
\label{b5:subsec:1301:02}

\begin{definition}\label{b5:def:invariant-bilinear-lie}
	设 \(\mathfrak g\) 是域 \(k\) 上的 Lie 代数，\(B:\mathfrak g\times\mathfrak g\rightarrow k\) 是双线性型。若
	\[
		B\bigl([x,y],z\bigr)=B\bigl(x,[y,z]\bigr)\qquad(x,y,z\in\mathfrak g),
	\]
	则 \(B\) 称为一个\term[bubian-shuangxianxing-xing]{不变双线性型}[invariant bilinear form]。
\end{definition}

\begin{proposition}\label{b5:prop:invariant-trace-form}
	设 \(\mathfrak g\) 为域 \(k\) 上的有限维 Lie 代数。它的每个有限维表示的迹型均为对称不变双线性型，因此其 Killing 型对称且不变。对任意对称不变双线性型 \(B:\mathfrak g\times\mathfrak g\to k\) 和理想 \(I\subseteq\mathfrak g\)，正交补
	\[
		I^\perp=\bigl\{\,x\in\mathfrak g\mid B(x,I)=0\,\bigr\}
	\]
	也是理想；特别地，\(B\) 的根 \(\mathfrak g^\perp\) 是理想。
\end{proposition}
\begin{proof}
	矩阵迹满足 \(\operatorname{tr}(AB)=\operatorname{tr}(BA)\)，并且
	\[
		\operatorname{tr}\bigl([A,B]C\bigr)=\operatorname{tr}\bigl(A[B,C]\bigr).
	\]
	对表示矩阵代入即得对称性与不变性。若 \(u\in I^\perp\)、\(x\in\mathfrak g\)、\(y\in I\)，则
	\[
		B\bigl([x,u],y\bigr)=-B\bigl(u,[x,y]\bigr)=0
	\]
	因为 \([x,y]\in I\)。于是 \([x,u]\in I^\perp\)。
\end{proof}

将\cref{b5:prop:lie-dual-tensor-hom}的对偶作用用于伴随表示，就得到 \(\mathfrak g^*\) 上的\term[yubansui-biaoshi]{余伴随表示}[coadjoint representation]：
\[
	(a\cdot\varphi)(y)=-\varphi([a,y]).
\]
下面把不变型所保持的信息写成这两个表示之间的同态。

\begin{proposition}\label{b5:prop:invariant-form-coadjoint}
	设 \(\mathfrak g\) 为域 \(k\) 上的有限维 Lie 代数，\(B:\mathfrak g\times\mathfrak g\to k\) 为双线性型。定义线性映射
	\[
		B^\flat:\mathfrak g\longrightarrow\mathfrak g^*, \qquad B^\flat(x)(y)=B(x,y).
	\]
	\symidx{\(B^\flat\)：双线性型诱导的对偶映射}
	则 \(B\) 不变，当且仅当 \(B^\flat\) 是从伴随表示到余伴随表示的同态。此时 \(\ker B^\flat=\{\,x\in\mathfrak g\mid B(x,\mathfrak g)=0\,\}\)，并且 \(B^\flat\) 为表示同构，当且仅当 \(B\) 非退化。
\end{proposition}
\begin{proof}
	交织条件在 \(y\) 处求值，恰为
	\[
		B([a,x],y)=-B(x,[a,y])\qquad(a,x,y\in\mathfrak g).
	\]
	利用 \([a,x]=-[x,a]\)，这等价于 \(B([x,a],y)=B(x,[a,y])\)，即定义中的不变性。这一步并未使用对称性。映射的核由 \(B(x,y)=0\) 对所有 \(y\) 成立的 \(x\) 组成，正是型的左根。最后，有限维性给出 \(\dim\mathfrak g=\dim\mathfrak g^*\)，所以左根为零恰好使 \(B^\flat\) 双射，此时型的右根也为零；双射表示同态的逆仍保持作用。
\end{proof}

\begin{proposition}\label{b5:prop:weight-orthogonality-invariant}
	设 \(\mathfrak g\) 是域 \(k\) 上的 Lie 代数，\(B:\mathfrak g\times\mathfrak g\to k\) 是对称不变双线性型，\(\mathfrak h\subseteq\mathfrak g\) 是子代数，\(\alpha,\beta\in\operatorname{Hom}_k(\mathfrak h,k)\)，\(x,y\in\mathfrak g\)。若
	\[
		[h,x]=\alpha(h)x,\qquad[h,y]=\beta(h)y \quad(h\in\mathfrak h)
	\]
	且 \(\alpha+\beta\ne0\)，则 \(B(x,y)=0\)。 若进一步 \(\mathfrak g\) 有限维、\(B\) 非退化，并且 \(\mathfrak g\) 是 \(\mathfrak h\) 的共同权空间
	\[
		\mathfrak g_\alpha= \{\,x\in\mathfrak g\mid[h,x]=\alpha(h)x\;\text{对所有}\;h\in\mathfrak h\,\}
	\]
	的直和，则 \(B\) 在 \(\mathfrak g_\alpha\times\mathfrak g_{-\alpha}\) 上的配对非退化，且这两个权空间维数相同。
\end{proposition}
\begin{proof}
	不变性给出
	\[
		0=B\bigl([h,x],y\bigr)+B\bigl(x,[h,y]\bigr) =\bigl(\alpha(h)+\beta(h)\bigr)B(x,y).
	\]
	取使 \(\alpha(h)+\beta(h)\ne0\) 的 \(h\)，即可除去该非零标量。

	在进一步的假设下，若 \(x\in\mathfrak g_\alpha\) 与 \(\mathfrak g_{-\alpha}\) 全部正交，前一结论又使它与其余权空间正交。任意向量都是这些权分量之和，故 \(x\) 与整个 \(\mathfrak g\) 正交，非退化性给出 \(x=0\)。交换两个权空间作同样论证，得到配对在两边均非退化。它们因而分别单射到对方的对偶，有限维性使两边维数相同。
\end{proof}

对偶作用把权 \(\beta\) 改成 \(-\beta\)，所以 \(B^\flat\) 的交织性质表现为相反权之间的配对；后面的根空间正交关系正要使用这个事实。在 \(\mathfrak{sl}_2(\mathbb C)\) 中，\([h,e]=2e\)、\([h,f]=-2f\)，所以 \(e\) 与 \(f\) 分属权 \(2\) 与 \(-2\) 的空间。前面的计算给出 \(\kappa(e,e)=\kappa(e,h)=0\)，而 \(\kappa(e,f)=4\)。向量 \(e\) 虽与自身配对为零，却不在型的根中；各向同性只要求 \(\kappa(e,e)=0\)，退化则要求某个非零向量与所有向量都配对为零。这两个条件不同。

\begin{lemma}\label{b5:lem:killing-restriction-ideal}
	设 \(\mathfrak g\) 是域 \(k\) 上的有限维 Lie 代数，\(I\) 是理想。对 \(x,y\in I\)，有
	\[
		\kappa_{\mathfrak g}(x,y)=\kappa_I(x,y).
	\]
\end{lemma}
\begin{proof}
	伴随算子 \(\operatorname{ad}x\) 把 \(\mathfrak g\) 送入 \(I\)，在 \(\mathfrak g/I\) 上诱导零作用。将 \(I\) 的基扩充为 \(\mathfrak g\) 的基，算子乘积的矩阵呈分块上三角形，商块为零，故其全空间迹等于限制到 \(I\) 的迹。
\end{proof}

\subsection{单 Lie 代数上的形式}
\label{b5:subsec:1301:03}

\begin{proposition}\label{b5:prop:simple-invariant-form-unique}
	设 \(\mathfrak g\) 是非零有限维复单 Lie 代数。则每个非零对称不变双线性型均非退化；所有对称不变双线性型组成的向量空间维数至多为一。
\end{proposition}
\begin{proof}
	由\cref{b5:prop:invariant-trace-form}，一个不变型的根是理想。非零型的根不可能是整个 \(\mathfrak g\)，单性便迫使其根为零。

	若对称不变双线性型只有零型，所成空间的维数为零。否则固定一个非零型 \(B_0\)。由\cref{b5:prop:invariant-form-coadjoint}，\(B_0^\flat\) 是伴随模到余伴随模的同构，而另一个型 \(B\) 给出的 \(B^\flat\) 是这两个模之间的同态。因此算子
	\[
		T=(B_0^\flat)^{-1}B^\flat
	\]
	是伴随模的自同态，具体满足 \(T[a,x]=[a,Tx]\) 及 \(B(x,y)=B_0(Tx,y)\)。由\cref{b5:prop:adjoint-representation}，伴随子模的条件是 \([\mathfrak g,U]\subseteq U\)，即 \(U\) 为理想；单性因此使这个非零伴随模不可约。这里使用的正是 Schur 引理在有限维复不可约模上的标量结论：复数域上 \(T\) 有特征值 \(c\)，非零空间 \(\ker(T-cI)\) 是伴随子模，故等于 \(\mathfrak g\)。于是 \(T=cI\)，\(B=cB_0\)。
\end{proof}

这时还不能仅凭此命题宣称 Killing 型非零：它只说明“若已有一个非零不变型，则唯一到标量”。下一节的判据将保证复单 Lie 代数确实具有非退化 Killing 型，从而把“至多一维”加强为“恰好一维”。

\ProseParagraphBreak

复数域在唯一性证明中保证自同态有特征值。在一般域上，Schur 引理只保证不可约模的自同态代数是除代数，未必只含底域标量。将 \(\mathfrak{sl}_2(\mathbb C)\) 看成实 Lie 代数 \(\mathfrak g_{\mathbb R}\)，就能看见这种差别。先核对它在实数域上仍单：对 \(A=\operatorname{ad}h\)，三个实特征空间为 \(\mathbb Ce,\mathbb Cf,\mathbb Ch\)，相应投影是
\[
	\dfrac{A(A+2\operatorname{id})}{8},\qquad \dfrac{A(A-2\operatorname{id})}{8},\qquad \operatorname{id}-\dfrac{A^2}{4}.
\]
任一非零实理想 \(J\) 被 \(A\) 及这些投影保持，因而含有其中一个空间的非零向量。若 \(ae\in J\)、\(a\in\mathbb C^\times\)，则 \([ih,ae]=2iae\) 使整个 \(\mathbb Ce\subseteq J\)，再由 \([f,e]=-h\)、\([h,f]=-2f\) 得到 \(\mathbb Ch,\mathbb Cf\subseteq J\)。\(af\in J\) 的情形相同；若 \(ah\in J\)，先用 \([ah,e]=2ae\) 回到第一种情形。故 \(J=\mathfrak g_{\mathbb R}\)，而 \([e,f]=h\ne0\)，所以它确为实单 Lie 代数。

\ProseParagraphBreak

记 \(\kappa_{\mathbb C}\) 为它原来的复 Killing 型，定义实值双线性型
\[
	B_1(x,y)=\operatorname{Re}\kappa_{\mathbb C}(x,y), \qquad B_2(x,y)=\operatorname{Im}\kappa_{\mathbb C}(x,y).
\]
取实部与虚部保留对称性及不变性，而
\[
	\begin{aligned}
		B_1(h,h)&=8,& B_2(h,h)&=0,\\
		B_1(ih,h)&=0,& B_2(ih,h)&=8
	\end{aligned}
\]
说明它们实线性无关。事实上两者都非退化：若 \(x\) 与所有 \(y\) 的实部配对为零，再把 \(y\) 换成 \(iy\)，虚部也为零，便由已算出的非退化 \(\kappa_{\mathbb C}\) 得 \(x=0\)；从虚部开始论证相同。因此实单性并不保证不变型只有一个实标量参数。这里乘以 \(i\) 的算子与全部伴随作用交换，却不是实标量，正是两种基域下自同态代数不同的具体表现。

% !TeX root = ../../Book5.tex
% 纲要标题：### 13.2 Cartan 判据
% 纲要位置：工作记录/计划纲要.md，第 4491 行。

\section{Cartan 判据}
\label{b5:sec:1302}

% 纲要范围：
%
% * 可解性的迹判据
% * 半单性的非退化判据
% * 与 Engel 定理的关系
% * 用谱与 Jordan 块计算说明高次插值，比较型的根与可解根、幂零根，并展开交换理想使 Killing 型退化的机制
%

迹型给出的二次条件足以判别可解性。不过，复数特征值的平方可能相消，不能从 \(\operatorname{tr}(x^2)=0\) 直接推出 \(x\) 幂零。证明中将另构造一个算子，把所用的迹写成特征值绝对值平方之和。

\subsection{可解性的迹判据}
\label{b5:subsec:1302:01}

\begin{lemma}[谱插值]\label{b5:lem:cartan-spectral-interpolation}
	设 \(V\) 为有限维复向量空间，\(x\in\operatorname{End}_{\mathbb C}(V)\)，\(V_\lambda\) 表示 \(x\) 对应于特征值 \(\lambda\) 的广义特征空间。按这些空间写 \(x=S+N\)，其中 \(S\) 在 \(V_\lambda\) 上为 \(\lambda I\)，\(N\) 幂零。令 \(S^\#\) 在 \(V_\lambda\) 上为 \(\overline\lambda I\)。则存在复多项式 \(p\)，满足 \(p(0)=0\)，并且
	\[
		\operatorname{ad}S^\#=p(\operatorname{ad}x) \quad\text{在}\;\operatorname{End}_{\mathbb C}(V)\;\text{上}.
	\]
\end{lemma}
\begin{proof}
	若 \(V=0\)，取 \(p=0\) 即可。以下设 \(V\ne0\)。 在 \(\operatorname{Hom}(V_\mu,V_\lambda)\) 上， \(\operatorname{ad}x=(\lambda-\mu)I+\operatorname{ad}N\)，其中后项由\cref{b5:lem:ad-of-nilpotent-matrix}幂零；而 \(\operatorname{ad}S^\#\) 是标量 \(\overline{\lambda-\mu}\)。取一个对所有这些子空间都足够大的整数 \(r\ge1\)。不同差值 \(\delta=\lambda-\mu\) 对应的 \((t-\delta)^r\) 两两互素，因此多项式中国剩余定理给出
	\[
		p(t)\equiv\overline\delta\pmod{(t-\delta)^r}.
	\]
	当 \(\delta=0\) 时，这要求 \(p(t)\equiv0\pmod{t^r}\)，特别有 \(p(0)=0\)。在相应 Hom 空间上记 \(Q=\operatorname{ad}N\)，则 \(Q^r=0\)。上述同余表示 \(p(t)-\overline\delta=(t-\delta)^r u_\delta(t)\)，所以
	\[
		p(\delta I+Q)-\overline\delta I =Q^r u_\delta(\delta I+Q)=0.
	\]
	因此 \(p(\operatorname{ad}x)\) 在整个广义特征块上都是所需标量。这些 Hom 空间的直和就是全部自同态空间，故算子等式成立。若只要求 \(p(\delta)=\overline\delta\)，一般还会留下 \(p'(\delta)Q\) 等项，不能得到这一结论。零差值处的高次同余同时保证 \(p(0)=0\)；下一定理将靠没有常数项这一性质，把 \(p(\operatorname{ad}x)(L)\) 留在导出理想中。
\end{proof}

例如取复空间 \(V=\mathbb C^2\) 上的算子
\[
	x=S=\operatorname{diag}(1,i),\qquad S^\#=\operatorname{diag}(1,-i).
\]
此时 \(x\) 不幂零，却有
\[
	\operatorname{tr}(x^2)=1+i^2=0,\qquad \operatorname{tr}(S^\#x)=|1|^2+|i|^2=2.
\]
共轭特征值的作用是消去平方相消的可能性。这里的 \(S^\#\) 保留原广义特征空间，只替换它们上的标量，仍是一个复线性算子；这个定义不使用矩阵的逐元素共轭，也不预先选定 Hermitian 内积后取伴随。具体地，取 \(P=\begin{psmallmatrix}1&i\\0&1\end{psmallmatrix}\)，令 \(x'=PxP^{-1}\)，则其半单部分 \(S'=x'\) 满足
\[
	x'=\begin{pmatrix}1&-1-i\\0&i\end{pmatrix},\qquad
	(S')^\#=PS^\#P^{-1}
	=\begin{pmatrix}1&1-i\\0&-i\end{pmatrix}.
\]
而 \(x'\) 的逐元素共轭矩阵的右上元是 \(-1+i\)，标准 Hermitian 内积下的共轭转置矩阵的右上元是零，都不同于 \((S')^\#\) 的右上元。

\ProseParagraphBreak

高次同余还要消去广义特征块中的幂零部分。取 \(V=\mathbb C^3\)，令
\[
	x=\begin{pmatrix}1&1&0\\0&1&0\\0&0&0\end{pmatrix},\qquad
	V_1=\mathbb Ce_1+\mathbb Ce_2,\quad V_0=\mathbb Ce_3.
\]
在以 \(E_{13},E_{23}\) 为基的 \(\operatorname{Hom}(V_0,V_1)\) 上，\(\operatorname{ad}x=I+Q\)，其中 \(Q=\begin{psmallmatrix}0&1\\0&0\end{psmallmatrix}\)、\(Q^2=0\)。差值为 \(1\)，所需的 \(\operatorname{ad}S^\#\) 为 \(I\)。普通插值选 \(p_0(t)=t\)，虽有 \(p_0(1)=1\)，却给出 \(p_0(I+Q)=I+Q\ne I\)。改取 \(q(t)=2t-t^2\)，就有
\[
	q(t)-1=-(t-1)^2,\qquad q(I+Q)=2(I+Q)-(I+Q)^2=I.
\]
这只处理了差值为 \(1\) 的这个块；完整的 \(p\) 必须同时满足其他差值处的同余。特别在零差值处，\(p(0)=0\) 还不够消去幂零部分，需要 \(p(t)\equiv0\pmod{t^r}\)。中国剩余定理把各块的这些高次要求同时实现，而 \(p(0)=0\) 又使下一证明中的算子多项式没有常数项。

\term[Cartan-kejie-panju]{Cartan 可解性判据}[Cartan's solvability criterion]如下。

\begin{theorem}[Cartan 可解性判据]\label{b5:thm:cartan-solvability}
	设 \(L\subseteq\mathfrak{gl}(V)\) 是有限维复向量空间 \(V\) 上的 Lie 子代数。则 \(L\) 可解，当且仅当
	\[
		\operatorname{tr}(xy)=0 \quad\bigl(x\in[L,L],\ y\in L\bigr).
	\]
	因而有限维复 Lie 代数 \(\mathfrak g\) 可解，当且仅当 \(\kappa_{\mathfrak g}\bigl([\mathfrak g,\mathfrak g],\mathfrak g\bigr)=0\)。
\end{theorem}
\begin{proof}
	若 \(V=0\)，则 \(L=0\)，结论成立。以下设 \(V\ne0\)。

	若 \(L\) 可解，\cref{b5:prop:solvable-derived-nilpotent}已经证明迹条件。

	反之，固定 \(x\in[L,L]\)，采用\cref{b5:lem:cartan-spectral-interpolation}的 \(S^\#\)。这个新算子未必属于 \(L\)，所以不能直接把它代入假设；我们利用它与 \(L\) 的括号应用迹条件。因为 \(\operatorname{ad}x\) 把 \(L\) 送入理想 \([L,L]\)，而 \(p\) 没有常数项，所以
	\[
		[S^\#,L]=p(\operatorname{ad}x)(L)\subseteq[L,L].
	\]
	把 \(x\) 写成有限和 \(\sum_j[a_j,b_j]\)，其中 \(a_j,b_j\in L\)。迹的循环性和假设给出
	\[
		\operatorname{tr}(S^\#x) =\sum_j\operatorname{tr}\bigl([S^\#,a_j]b_j\bigr)=0.
	\]
	另一方面，写 \(x=S+N\)。在每个广义特征空间 \(V_\lambda\) 上，\(S=\lambda I\)、\(S^\#=\overline\lambda I\)，而 \(N|_{V_\lambda}\) 幂零，故其迹为零。因此
	\[
		\begin{aligned}
			\operatorname{tr}(S^\#N) &=\sum_\lambda\overline\lambda\operatorname{tr}(N|_{V_\lambda})=0,\\
			\operatorname{tr}(S^\#x) &=\operatorname{tr}(S^\#S)+\operatorname{tr}(S^\#N) =\sum_\lambda(\dim V_\lambda)|\lambda|^2.
		\end{aligned}
	\]
	每一项均非负，故每个特征值都为零，\(x\) 幂零。于是 Lie 子代数 \([L,L]\) 中的每个矩阵都是幂零算子。\cref{b5:thm:engel-common-zero}给出的完全旗标使其同时严格上三角化，故它是幂零 Lie 代数，特别可解。\([L,L]\) 是 \(L\) 的理想，商 \(L/[L,L]\) 交换，因而也可解；由\cref{b5:prop:lie-solvable-extension}，\(L\) 可解。

	最后，对抽象 Lie 代数取 \(L=\operatorname{ad}\mathfrak g\)，矩阵迹条件正好是陈述中的 Killing 型条件。若该条件成立，已经证明 \(L\) 可解，而 \(\mathfrak g/Z(\mathfrak g)\cong L\)。中心交换，由\cref{b5:prop:lie-solvable-extension}，\(\mathfrak g\) 可解。逆向对伴随表示使用已证的必要性。
\end{proof}

\subsection{半单性的非退化判据}
\label{b5:subsec:1302:02}

\term[Cartan-bandan-panju]{Cartan 半单性判据}[Cartan's semisimplicity criterion]给出一个更简洁的检验。

\begin{theorem}[Killing 型的非退化判据]\label{b5:thm:killing-semisimple}
	有限维复 Lie 代数 \(\mathfrak g\) 半单，当且仅当其 Killing 型非退化。
\end{theorem}
\begin{proof}
	设 \(\mathfrak g\) 半单，令 \(R=\mathfrak g^\perp\) 为 Killing 型的根。\cref{b5:prop:invariant-trace-form}说明 \(R\) 是理想。\cref{b5:lem:killing-restriction-ideal}给出 \(\kappa_R=\kappa_{\mathfrak g}|_{R\times R}=0\)，故\cref{b5:thm:cartan-solvability}说明 \(R\) 可解。半单性迫使 \(R=0\)。

	反之，设 Killing 型非退化。若有非零可解理想 \(I\)，其导出列的最后一个非零项 \(A\) 是非零交换理想：各导出项均为 \(\mathfrak g\) 的理想，且 \([A,A]=0\)。固定 \(a\in A\)、\(x\in\mathfrak g\)，令 \(T=\operatorname{ad}a\operatorname{ad}x\)。对任意 \(y\in\mathfrak g\)，理想条件给出
	\[
		T(y)=[a,[x,y]]\in A.
	\]
	若 \(y\in A\)，则 \([x,y]\in A\)，再用交换性便有 \([a,[x,y]]=0\)。因此 \(T(\mathfrak g)\subseteq A\)、\(T(A)=0\)，所以 \(T^2=0\)，\(\kappa_{\mathfrak g}(a,x)=\operatorname{tr}T=0\)。这使 \(A\) 包含于型的根，与非退化性矛盾。因此不存在非零可解理想。
\end{proof}

这里型的根 \(R=\{\,a\in\mathfrak g\mid\kappa_{\mathfrak g}(a,\mathfrak g)=0\,\}\) 由迹配对定义；\(\operatorname{rad}\mathfrak g\) 与 \(\operatorname{nilrad}\mathfrak g\) 则分别由最大可解理想、最大幂零理想定义。上面关于 \(R\) 的可解性论证对任意有限维复 Lie 代数都成立，因而只给出 \(R\subseteq\operatorname{rad}\mathfrak g\)，并不说明它们相等。一般表示的迹型还可能有另一种根。对 \(\mathfrak{gl}_n(\mathbb C)\)（\(n\ge1\)），\cref{b5:exm:killing-sln}之后已说明自然迹型 \(\beta(X,Y)=\operatorname{tr}(XY)\) 的根为零，但非零中心 \(\mathbb CI\) 是交换理想，所以包含于幂零根与可解根。一般迹型的根取决于所选表示，而后两个理想是 Lie 代数自身的结构。

\begin{example}\label{b5:exm:killing-semidirect-natural}
	在复向量空间 \(\mathfrak g=\mathfrak{sl}_2(\mathbb C)\oplus\mathbb C^2\) 上规定
	\[
		[(X,v),(Y,w)]=([X,Y],Xw-Yv).
	\]
	矩阵作用满足 \([X,Y]v=X(Yv)-Y(Xv)\)，结合 \(\mathfrak{sl}_2\) 的 Jacobi 恒等式，便保证上述括号也满足 Jacobi 恒等式。这是自然作用给出的半直积 Lie 代数，\(A=0\oplus\mathbb C^2\) 是交换理想。对 \(v\in\mathbb C^2\)，记 \(B_v:\mathfrak{sl}_2\to\mathbb C^2\)、\(B_v(Y)=-Yv\)。相对于上述向量空间分解，伴随算子为
	\[
		\operatorname{ad}(X,v)=
		\begin{pmatrix}\operatorname{ad}_{\mathfrak{sl}_2}X&0\\B_v&X\end{pmatrix}.
	\]
	若 \(a=(0,u)\in A\)，则对任意 \(x=(X,v)\in\mathfrak g\)，有
	\[
		\operatorname{ad}a\operatorname{ad}x
		=\begin{pmatrix}0&0\\B_u\operatorname{ad}_{\mathfrak{sl}_2}X&0\end{pmatrix}.
	\]
	这个矩阵把全空间送入 \(A\)，又零化 \(A\)，故平方为零、迹为零。这不是要求每个这样的乘积本身为零。例如取标准基 \(u_1,u_2\) 及 \(h=\operatorname{diag}(1,-1)\)、\(e=E_{12}\)、\(f=E_{21}\)，在基 \(h,e,f,u_1,u_2\) 下，有
	\[
		B_{u_1}=\begin{pmatrix}-1&0&0\\0&0&-1\end{pmatrix},\qquad
		\operatorname{ad}(h,0)=\operatorname{diag}(0,2,-2,1,-1).
	\]
	于是 \(\operatorname{ad}(0,u_1)\operatorname{ad}(h,0)\) 将 \((f,0)\) 送到 \((0,2u_2)\ne0\)，但仍具有上面的零对角块，因而迹为零。对所有 \(a\in A\)、\(x\in\mathfrak g\) 均得 \(\kappa_{\mathfrak g}(a,x)=0\)，所以 \(A\) 包含于 Killing 型的根。商 \(\mathfrak g/A\cong\mathfrak{sl}_2\)；它的单性将在\cref{b5:prop:classical-simple}中由矩阵单位证明。商的单性并不能消除原代数中这个交换理想造成的退化。
\end{example}

\begin{corollary}\label{b5:cor:simple-killing-forms}
	非零有限维复单 Lie 代数 \(\mathfrak g\) 的对称不变双线性型均为 Killing 型的标量倍数，且 Killing 型非退化。
\end{corollary}
\begin{proof}
	单 Lie 代数半单，所以\cref{b5:thm:killing-semisimple}保证 Killing 型非退化；再用\cref{b5:prop:simple-invariant-form-unique}。
\end{proof}

\begin{proposition}\label{b5:prop:faithful-trace-nondegenerate}
	设 \(\mathfrak g\) 是有限维复半单 Lie 代数，\(\rho:\mathfrak g\rightarrow\mathfrak{gl}(V)\) 是忠实有限维表示。则 \(\beta_\rho(x,y)=\operatorname{tr}\bigl(\rho(x)\rho(y)\bigr)\) 非退化。
\end{proposition}
\begin{proof}
	型的根 \(R\) 是理想，且 \(\operatorname{tr}\bigl(\rho(x)\rho(y)\bigr)=0\) 对全部 \(x,y\in R\) 成立。对矩阵代数 \(\rho(R)\) 用\cref{b5:thm:cartan-solvability}，得 \(\rho(R)\) 可解。忠实性给出 \(R\cong\rho(R)\)，故 \(R\) 是可解理想，只能为零。
\end{proof}

\subsection{与 Engel 定理的关系}
\label{b5:subsec:1302:03}

证明的顺序值得区分。Engel 定理先由共同零向量识别幂零性；Lie 定理随后给出可解表示的上三角形；Cartan 判据再通过谱插值把迹条件还原为 Engel 定理的假设。半单性判据用的是这条已经建立的证明链，并未调用半单 Lie 代数表示的完全可约性。

\begin{example}\label{b5:exm:killing-affine-two}
	设 \(\mathfrak g=\mathbb Ch+\mathbb Ce\)，\([h,e]=e\)。相对于 \(h,e\)，
	\[
		\operatorname{ad}h=\begin{pmatrix}0&0\\0&1\end{pmatrix},
		\qquad
		\operatorname{ad}e=\begin{pmatrix}0&0\\-1&0\end{pmatrix},
		\qquad
		[\kappa]_{h,e}=\begin{pmatrix}1&0\\0&0\end{pmatrix}.
	\]
	它可解而 Killing 型不恒为零；正确条件是 \(\kappa\bigl([\mathfrak g,\mathfrak g],\mathfrak g\bigr)=0\)，这里 \([\mathfrak g,\mathfrak g]=\mathbb Ce\) 恰好等于型的根。由\cref{b5:exm:nilradical-quotient-failure}中对同一代数的计算，\(\operatorname{rad}\mathfrak g=\mathfrak g\)、\(\operatorname{nilrad}\mathfrak g=\mathbb Ce\)，所以
	\[
		\mathfrak g^\perp=\operatorname{nilrad}\mathfrak g =\mathbb Ce\subsetneq\operatorname{rad}\mathfrak g=\mathfrak g.
	\]
	这里型的根严格小于可解根；它与幂零根相等只是这个例子的具体结果。
\end{example}

\begin{example}\label{b5:exm:zero-killing-nonnilpotent}
	设复 Lie 代数 \(\mathfrak g\) 有基 \(h,x,y\)，括号由
	\[
		[h,x]=x,\qquad[h,y]=iy,\qquad[x,y]=0
	\]
	给出。在此基下，伴随算子为
	\[
		\operatorname{ad}h=\operatorname{diag}(0,1,i),\qquad \operatorname{ad}x=-E_{21},\qquad \operatorname{ad}y=-iE_{31}.
	\]
	后两个矩阵与任意伴随矩阵相乘的迹均为零，而 \(\operatorname{tr}((\operatorname{ad}h)^2)=1+i^2=0\)。由双线性性，Killing 型恒为零，所以其根为整个 \(\mathfrak g\)。导出代数 \(A=\mathbb Cx+\mathbb Cy\) 交换，故 \(\mathfrak g\) 可解，\(\operatorname{rad}\mathfrak g=\mathfrak g\)。然而 \((\operatorname{ad}h)^m x=x\ne0\) 对任意 \(m\ge1\) 成立，因而 \(\mathfrak g\) 不幂零。

	理想 \(A\) 交换，所以 \(A\subseteq\operatorname{nilrad}\mathfrak g\)。若 \(z=ah+bx+cy\) 属于一个幂零理想，\cref{b5:prop:lie-radicals-exist}证明中的伴随作用计算说明 \(\operatorname{ad}z\) 在整个 \(\mathfrak g\) 上幂零。可是 \([z,x]=ax\)，幂零性迫使 \(a=0\)。因此每个幂零理想都包含于 \(A\)，得到
	\[
		\operatorname{nilrad}\mathfrak g=A\subsetneq\mathfrak g =\mathfrak g^\perp=\operatorname{rad}\mathfrak g.
	\]
	这个例子的 Killing 根严格大于幂零根，故前一例中的相等关系不能一般化。
\end{example}

特征零假设同样不可漏写。以特征 \(3\) 的域 \(k\) 为例，令 \(e=E_{12}\)、\(h=E_{11}-E_{22}\)、\(f=E_{21}\)，它们构成 \(\mathfrak{sl}_2(k)\) 的一组基。\(\operatorname{ad}h\) 在这三个基向量上的特征值为 \(2,0,-2=1\)，仍互不相同。对 \(\operatorname{ad}h\) 作 Lagrange 插值，可将各特征空间的投影写成它的多项式。任意理想被 \(\operatorname{ad}h\) 保持，因而也被这些投影保持；对非零理想中的一个非零元素投影，至少得到一个非零基向量的标量倍数。再由 \([h,e]=2e\)、\([h,f]=-2f\)、\([e,f]=h\) 得到全部三个基向量，所以 \(\mathfrak{sl}_2(k)\) 仍是单 Lie 代数。但若特征为 \(2\)，这些步骤与 Killing 型公式都会退化，\(\mathfrak{sl}_2\) 不再半单。因此本章判据的适用范围始终是有限维复 Lie 代数，不能只看相同的矩阵记号就推广。

% !TeX root = ../../Book5.tex
% 纲要标题：### 13.3 半单分解
% 纲要位置：工作记录/计划纲要.md，第 4498 行。

\section{半单分解}
\label{b5:sec:1303}

% 纲要范围：
%
% * 半单 Lie 代数分解为单理想
% * 理想与正交补
% * 分解的唯一性
% * 用单理想的正交配对分类不变型，并由半单商确定一般表示迹型的根
%

一般表示中的不变子空间未必有不变补。伴随表示在半单情形却有一个特别直接的补：用 Killing 型取正交补。这是半单结构分解的来源，而非任意表示完全可约性的预先使用。

\subsection{半单 Lie 代数分解为单理想}
\label{b5:subsec:1303:01}

\begin{theorem}[半单 Lie 代数的理想分解]\label{b5:thm:semisimple-ideal-decomposition}
设 \(\mathfrak g\) 是有限维复半单 Lie 代数。则存在单理想
\(\mathfrak g_1,\ldots,\mathfrak g_r\)，使
\[
\mathfrak g=\mathfrak g_1\oplus\cdots\oplus\mathfrak g_r,
\qquad[\mathfrak g_i,\mathfrak g_j]=0\quad(i\ne j).
\]
此外，\(\mathfrak g\) 的每个理想都是其中若干 \(\mathfrak g_i\) 的直和。这些单理想作为 \(\mathfrak g\) 的子空间由 \(\mathfrak g\) 唯一确定，只允许改变排列顺序。
\end{theorem}
\begin{proof}
先取任意理想 \(I\)。由\cref{b5:prop:invariant-trace-form}，\(I^\perp\) 是理想，所以 \(J=I\cap I^\perp\) 也是理想。这里 \(J\) 是限制型 \(\kappa_{\mathfrak g}|_{I\times I}\) 的根；要排除它，仍须证明它是可解理想。对 \(x,y\in J\)，有 \(\kappa_{\mathfrak g}(x,y)=0\)，再由\cref{b5:lem:killing-restriction-ideal}得到 \(\kappa_J=0\)。Cartan 可解性判据说明 \(J\) 可解，半单性迫使 \(J=0\)。由于整个 Killing 型非退化，
\(\dim I^\perp=\dim\mathfrak g-\dim I\)，故
\[
\mathfrak g=I\oplus I^\perp.
\]
两个理想的括号同时属于二者，因而为零。

如果 \(\mathfrak g\ne0\)，选一个极小非零理想 \(I\)。若 \(K\) 是 \(I\) 的 Lie 理想，则 \([I,K]\subseteq K\)，而 \([I^\perp,K]=0\)。由 \(\mathfrak g=I\oplus I^\perp\) 得 \([\mathfrak g,K]\subseteq K\)，所以 \(K\) 也是 \(\mathfrak g\) 的理想。极小性说明 \(I\) 没有非零真理想；它不能交换，否则会是可解理想。因此 \(I\) 为单 Lie 代数。同理，补 \(I^\perp\) 的每个可解理想也是 \(\mathfrak g\) 的可解理想，故补仍半单。对补继续此过程，维数严格减小，有限步便得到所述分解。

设 \(J\) 是任意理想，\(x=x_1+\cdots+x_r\in J\)。若 \(x_i\ne0\)，单 Lie 代数 \(\mathfrak g_i\) 的中心为零，所以
\([\mathfrak g_i,x_i]\ne0\)。又因为不同分量交换，
\[
[\mathfrak g_i,x]=[\mathfrak g_i,x_i]\subseteq J\cap\mathfrak g_i.
\]
交集是 \(\mathfrak g_i\) 的理想，故单性迫使 \(\mathfrak g_i\subseteq J\)。于是 \(x\) 的每个非零分量都属于已包含于 \(J\) 的单理想，从而 \(J\) 就是这些理想的直和。特别地，极小非零理想恰好是所列的 \(\mathfrak g_i\)，唯一性也随之成立。
\end{proof}

\subsection{理想与正交补}
\label{b5:subsec:1303:02}

\begin{corollary}\label{b5:cor:semisimple-ideal-complement}
设 \(\mathfrak g\) 是有限维复半单 Lie 代数，\(I\) 是理想。则
\[
\mathfrak g=I\oplus I^\perp,\qquad
I^\perp=Z_{\mathfrak g}(I),\qquad
\mathfrak g/I\cong I^\perp.
\]
这里正交补取自 Killing 型，中心化子为
\(Z_{\mathfrak g}(I)=\bigl\{\,x\mid[x,I]=0\,\bigr\}\)。
而且 \(I\) 与 \(\mathfrak g/I\) 都半单。
\end{corollary}
\begin{proof}
由\cref{b5:thm:semisimple-ideal-decomposition}，\(I\) 是部分单理想的和，其余分量组成 \(I^\perp\)。后者显然中心化 \(I\)；若一个元素在 \(I\) 分量上的投影也中心化 \(I\)，该投影属于 \(Z(I)=0\)，所以中心化子没有更多元素。商映射在 \(I^\perp\) 上单射且满射，给出同构。两者都由单理想直和组成，故半单。
\end{proof}

\begin{corollary}\label{b5:cor:semisimple-perfect}
有限维复半单 Lie 代数 \(\mathfrak g\) 满足
\[
[\mathfrak g,\mathfrak g]=\mathfrak g,\qquad Z(\mathfrak g)=0.
\]
因此其一维表示全为平凡表示。
\end{corollary}
\begin{proof}
每个非交换单理想的导出代数是非零理想，因而等于自身；中心是可解理想，因而为零。对\cref{b5:thm:semisimple-ideal-decomposition}中的直和逐项应用即可。一维表示的像交换，必零化导出代数，故零化整个 \(\mathfrak g\)。
\end{proof}

分解以后，其他对称不变型也可以逐个单理想确定。关键不只是每个单因子上的型唯一，还要说明不同因子之间不能有交叉配对；即使两个因子同构，不变性也会排除这种配对。

\begin{proposition}\label{b5:prop:semisimple-invariant-forms}
设有限维复半单 Lie 代数的单理想分解为
\(\mathfrak g=\bigoplus_{i=1}^r\mathfrak g_i\)。其全部对称不变双线性型恰为
\[
B=\bigoplus_{i=1}^r c_i\kappa_{\mathfrak g_i},\qquad c_i\in\mathbb C.
\]
该型的根为 \(\bigoplus_{c_i=0}\mathfrak g_i\)，因此非退化当且仅当每个 \(c_i\ne0\)。
\end{proposition}
\begin{proof}
若 \(i\ne j\)，取 \(a,b\in\mathfrak g_i\) 和 \(y\in\mathfrak g_j\)。不同理想交换，不变性给出
\[
B\bigl([a,b],y\bigr)=B\bigl(a,[b,y]\bigr)=0.
\]
而 \([\mathfrak g_i,\mathfrak g_i]=\mathfrak g_i\)，即每个 \(x\in\mathfrak g_i\) 都是有限个括号的和，所以 \(B(x,y)=0\)。这一步只使用不同理想相互交换，以及 \([\mathfrak g_i,\mathfrak g_i]=\mathfrak g_i\) 这一性质（通常称为完美性），与两个因子是否同构无关。于是 \(B\) 是各限制型的正交直和；由\cref{b5:cor:simple-killing-forms}，每个限制恰为 \(c_i\kappa_{\mathfrak g_i}\)。反过来，这些型的任意直和都对称且不变。各单因子的 Killing 型非退化，故只有系数为零的分量进入型的根，得到最后两项结论。
\end{proof}

对表示迹型，系数为零的分量恰好是表示完全看不见的单理想。忠实性因此有一个具体作用：它排除了这些消失的分量。

\begin{corollary}\label{b5:cor:semisimple-trace-kernel}
设 \(\mathfrak g\) 是有限维复半单 Lie 代数，\(\rho\) 是其任意有限维表示。则
\[
\operatorname{rad}(\beta_\rho)=\ker\rho.
\]
\end{corollary}
\begin{proof}
记 \(K=\ker\rho\)。由\cref{b5:cor:semisimple-ideal-complement}，商 \(\mathfrak g/K\) 仍半单。表示在商上诱导忠实表示
\(\overline\rho(x+K)=\rho(x)\)，其迹型由\cref{b5:prop:faithful-trace-nondegenerate}非退化。原迹型正是沿商映射拉回的型。具体地，\(x\) 与所有 \(y\in\mathfrak g\) 的配对为零，当且仅当 \(x+K\) 在商上与全部元素正交；商上的非退化性迫使 \(x+K=0\)，即 \(x\in K\)。反向包含也由 \(\rho(x)=0\) 立即成立。
\end{proof}

例如，设 \(\mathfrak s\) 是非零有限维复单 Lie 代数，在
\(\mathfrak g=\mathfrak s\oplus\mathfrak s\) 上取作用于 \(V=\mathfrak s\) 的表示
\(\rho(x,y)=\operatorname{ad}_{\mathfrak s}x\)。则
\[
\beta_\rho\bigl((x,y),(u,v)\bigr)=\kappa_{\mathfrak s}(x,u),\qquad
\ker\rho=0\oplus\mathfrak s.
\]
它只检测第一因子，故迹型在第二因子上恒零；环境代数半单并不能使这个非忠实表示的迹型非退化。

理想的 Killing 正交补则总能给出上述直和分解。一般子代数 \(\mathfrak h\) 的正交补不一定是理想，\(\kappa|_{\mathfrak h}\) 也不一定非退化；例如 \(\mathfrak{sl}_2\) 中的一维子代数 \(\mathbb Ce\) 就是各向同性的。

\subsection{分解的唯一性}
\label{b5:subsec:1303:03}

\begin{example}\label{b5:exm:diagonal-simple-not-ideal}
设 \(\mathfrak s\) 是非零复单 Lie 代数。对角子代数
\[
\Delta\mathfrak s=\bigl\{\,(x,x)\mid x\in\mathfrak s\,\bigr\}
\subseteq\mathfrak s\oplus\mathfrak s
\]
同构于 \(\mathfrak s\)，但不是理想。取 \([a,x]\ne0\)，则
\[
\bigl[(a,0),(x,x)\bigr]=\bigl([a,x],0\bigr)\notin\Delta\mathfrak s.
\]
所以单子代数可以很多，而理想分解中的单理想仍只有两个。
\end{example}

分解唯一性说的是具体理想的集合唯一，比仅仅列出抽象同构类型更强。自同构可以置换同构的单理想，但不可能把一个单理想送到上述对角子代数，因为自同构保持理想性。这也说明分解中出现重数时仍没有连续选择补空间的自由。

% !TeX root = ../../Book5.tex
% 纲要标题：### 13.4 经典例子
% 纲要位置：工作记录/计划纲要.md，第 4505 行。

\section{经典例子}
\label{b5:sec:1304}

% 纲要范围：
%
% * sl_n、so_n 与 sp_2n
% * 维数、中心与单性的范围
% * 小秩同构及例外情形
% * 用括号传播与特征标投影解释单性证明，并以外幂配对展开四维低秩同构
%
% ---
%

经典矩阵 Lie 代数把上述判据变成可以直接计算的例子。维数、中心和单性可以由矩阵单位、型及伴随作用直接证明。后面的根系将解释这些计算为什么分成四个系列。

\subsection{\texorpdfstring{\(\mathfrak{sl}_n\)、\(\mathfrak{so}_n\) 与 \(\mathfrak{sp}_{2n}\)}{sln、son 与 sp2n}}
\label{b5:subsec:1304:01}

对复向量空间，非退化对称型可选基写成单位矩阵，非退化交替型可写成标准辛矩阵。因此采用下列矩阵实现：
\[
\begin{aligned}
\mathfrak{sl}_n(\mathbb C)&=\bigl\{\,X\in M_n(\mathbb C)\mid\operatorname{tr}X=0\,\bigr\},\\
\mathfrak{so}_n(\mathbb C)&=\bigl\{\,X\in M_n(\mathbb C)\mid X^{\mathsf T}+X=0\,\bigr\},\\
\mathfrak{sp}_{2n}(\mathbb C)&=\bigl\{\,X\in M_{2n}(\mathbb C)\mid
X^{\mathsf T}J+JX=0\,\bigr\},\qquad
J=\begin{pmatrix}0&I_n\\-I_n&0\end{pmatrix}.
\end{aligned}
\]
换用另一个非退化型，只会得到由换基共轭而同构的 Lie 代数。

\ProseParagraphBreak

这里保持的是复双线性型，因此使用普通转置 \(X^{\mathsf T}\)。若保持标准 Hermitian 型，相应条件则是 \(X^*+X=0\)，其中 \(X^*\) 为共轭转置。例如对 \(n\ge2\)，令 \(F=E_{12}-E_{21}\)，则 \(F^{\mathsf T}=-F\)，而 \((iF)^{\mathsf T}=-iF\)，所以 \(F,iF\) 都属于 \(\mathfrak{so}_n(\mathbb C)\)。但 \((iF)^*=iF\)，故 \(iF\) 不是斜 Hermitian 矩阵。前者对复标量封闭，后者只构成实 Lie 代数；尤其不能把这里的复 \(\mathfrak{so}_3\) 与实旋转 Lie 代数 \(\mathfrak{so}_3(\mathbb R)\) 混同。

\begin{proposition}\label{b5:prop:classical-dimensions}
对整数 \(n\ge1\)，复经典 Lie 代数 \(\mathfrak{sl}_n(\mathbb C)\)、\(\mathfrak{so}_n(\mathbb C)\)、\(\mathfrak{sp}_{2n}(\mathbb C)\) 的维数为
\[
\dim\mathfrak{sl}_n=n^2-1,\qquad
\dim\mathfrak{so}_n=\frac{n(n-1)}2,\qquad
\dim\mathfrak{sp}_{2n}=n(2n+1).
\]
相对于标准辛型，\(\mathfrak{sp}_{2n}(\mathbb C)\) 恰好由如下分块矩阵构成，其中 \(A,B,C\in M_n(\mathbb C)\)：
\[
X=\begin{pmatrix}A&B\\C&-A^{\mathsf T}\end{pmatrix},
\qquad B=B^{\mathsf T},\quad C=C^{\mathsf T}.
\]
\end{proposition}
\begin{proof}
迹零是 \(M_n\) 上一个非零线性条件，所以第一式成立。反对称矩阵的对角元为零，每对 \(i<j\) 留一个自由参数，得到第二式。把四个矩阵块代入 \(X^{\mathsf T}J+JX=0\)，便得陈述中的块条件；自由参数数目为
\(n^2+2\cdot n(n+1)/2=n(2n+1)\)。
\end{proof}

\subsection{维数、中心与单性的范围}
\label{b5:subsec:1304:02}

\begin{lemma}[提取共同权分量]\label{b5:lem:classical-weight-projection}
设有限维复向量空间 \(W\) 上的有限族算子 \(H_1,\ldots,H_r\) 交换且可同时对角化。若子空间 \(U\subseteq W\) 被它们保持，则每个 \(u\in U\) 的各个共同特征空间分量仍属于 \(U\)。
\end{lemma}
\begin{proof}
共同特征值元组只有有限多个。可选系数 \(c_i\)，使不同元组在
\(H=\sum_i c_iH_i\) 下具有不同特征值：需要排除的只是有限多个真超平面。对每个特征值 \(\lambda\)，Lagrange 插值多项式
\[
P_\lambda(t)=\prod_{\mu\ne\lambda}\frac{t-\mu}{\lambda-\mu}
\]
给出相应投影 \(P_\lambda(H)\)。子空间 \(U\) 被 \(H\) 及其多项式保持，故被这些投影保持。
\end{proof}

\begin{proposition}\label{b5:prop:classical-simple}
在复数域上，\(\mathfrak{sl}_n\) 对 \(n\ge2\) 为单 Lie 代数，
\(\mathfrak{sp}_{2n}\) 对 \(n\ge1\) 为单 Lie 代数，
\(\mathfrak{so}_n\) 对 \(n=3\) 或 \(n\ge5\) 为单 Lie 代数。因此这些范围内的中心均为零。
\end{proposition}
\begin{proof}
先考虑特殊线性 Lie 代数。
令 \(\mathfrak h\) 为迹零对角矩阵空间。对 \(H=\operatorname{diag}(t_1,\ldots,t_n)\)，有
\[
[H,E_{ij}]=(t_i-t_j)E_{ij}.
\]
这些非零权彼此不同，零权空间为 \(\mathfrak h\)。给定非零理想 \(I\)，\cref{b5:lem:classical-weight-projection}说明：若 \(I\) 的元素有非零非对角分量，则某个 \(E_{ij}\in I\)；否则 \(I\) 含有非零迹零对角矩阵 \(H\)，其对角元不能全相等，故 \([H,E_{ij}]\ne0\) 也使 \(E_{ij}\in I\)。

令 \(H_{ij}=E_{ii}-E_{jj}\)。由 \([E_{ij},E_{ji}]=H_{ij}\) 得到 \(H_{ij}\in I\)，再由 \([H_{ij},E_{ji}]=-2E_{ji}\) 及特征零，得 \(E_{ji}\in I\)。对 \(k\ne i,j\)，与 \(H_{ij}\) 括号先给出 \(E_{ik},E_{ki},E_{jk},E_{kj}\in I\)；随后
\([E_{ki},E_{i\ell}]=E_{k\ell}\) 给出其余非对角单位。全部迹零对角矩阵由 \([E_{ij},E_{ji}]\) 生成，所以 \(I=\mathfrak{sl}_n\)。当 \(n=2\) 时没有额外的 \(k\)，前三个矩阵已经构成一组基。

再考虑辛 Lie 代数。
用 \(\bar i=n+i\) 表示后半部指标，记
\[
\begin{aligned}
H_i&=E_{ii}-E_{\bar i\bar i},&
A_{ij}&=E_{ij}-E_{\bar j\bar i}\quad(i\ne j),\\
B_{ii}&=E_{i\bar i},&
C_{ii}&=E_{\bar i i},\\
B_{ij}&=E_{i\bar j}+E_{j\bar i},&
C_{ij}&=E_{\bar i j}+E_{\bar j i}\quad(i<j).
\end{aligned}
\]
这些矩阵构成一组基。相对于 \(H=\sum_i t_iH_i\) 的伴随作用，它们的权分别是
\(0,t_i-t_j,2t_i,-2t_i,t_i+t_j,-t_i-t_j\)，所有非零权空间均为一维。给定非零理想 \(I\)，若其某个元素有非零对角分量，投影引理把这个零权分量提取为 \(I\) 中的非零对角元；否则，先提取其中一个非零权向量。各非零权空间为一维，所以除以非零系数后，可在 \(I\) 中取得相应的基权向量。再与整个代数中相反权的基向量取括号，理想性质保证结果仍在 \(I\) 中。相应的三种配对为
\[
[A_{ij},A_{ji}]=H_i-H_j,\quad
[B_{ij},C_{ij}]=H_i+H_j,\quad
[B_{ii},C_{ii}]=H_i,
\]
所以无论如何，\(I\) 总含有一个非零对角元 \(H=\sum_i t_iH_i\)。

选择 \(t_i\ne0\)，由 \([H,B_{ii}]=2t_iB_{ii}\)、
\([H,C_{ii}]=-2t_iC_{ii}\) 得 \(B_{ii},C_{ii},H_i\in I\)。
当 \(n=1\) 时，\(H_1,B_{11},C_{11}\) 已是一组基，故 \(I=\mathfrak{sp}_2\)。以下设 \(n\ge2\)。
对 \(j\ne i\)，与 \(H_i\) 括号得到 \(A_{ij},A_{ji}\in I\)，再得到 \(H_i-H_j\)，故所有 \(H_j\) 均在 \(I\) 中。每个非零权至少在某个 \(H_j\) 上非零，与这个 \(H_j\) 括号便把相应基向量收入 \(I\)。所以 \(I=\mathfrak{sp}_{2n}\)。

在 \(\mathfrak{sp}_4\) 中，可以从一个具体的非零权向量 \(B_{11}\in I\) 看见这次传播。矩阵单位的乘法首先给出
\[
\begin{aligned}
[B_{11},C_{11}]&=H_1,&[H_1,C_{11}]&=-2C_{11},\\
[H_1,A_{12}]&=A_{12},&[H_1,A_{21}]&=-A_{21}.
\end{aligned}
\]
因此 \(H_1,C_{11},A_{12},A_{21}\in I\)。再由
\[
[A_{12},A_{21}]=H_1-H_2
\]
得到 \(H_2\in I\)。最后，
\[
\begin{aligned}
[H_2,B_{22}]&=2B_{22},&[H_2,C_{22}]&=-2C_{22},\\
[H_1,B_{12}]&=B_{12},&[H_1,C_{12}]&=-C_{12}
\end{aligned}
\]
把剩下的四个向量也收入 \(I\)。连同最初的 \(B_{11}\)，这正是全部十个基向量。每一步只需括号的一端已在 \(I\) 中，另一端可以是整个代数中的任意元素，这也是使用理想条件的地方。

最后考虑正交 Lie 代数。
记 \(F_{ij}=E_{ij}-E_{ji}\)，其中 \(i<j\)。考虑阶为 \(2^{n-1}\) 的有限交换群
\[
D=\left\{\,\operatorname{diag}(\varepsilon_1,\ldots,\varepsilon_n)
\mid\varepsilon_i=\pm1,\ \prod_i\varepsilon_i=1\,\right\}.
\]
共轭作用将 \(F_{ij}\) 乘以 \(\varepsilon_i\varepsilon_j\)。每个 \(D\) 的元素都是若干同时翻转两个坐标符号的矩阵之积。记同时翻转第 \(i,j\) 个坐标的矩阵为 \(d_{ij}\)，则对 \(X\in\mathfrak{so}_n\) 有
\[
d_{ij}Xd_{ij}^{-1}=X+2[F_{ij},[F_{ij},X]].
\]
这个等式可在基向量上验证：\((\operatorname{ad}F_{ij})^2\) 在只含 \(i,j\) 中一个指标的基向量上为负恒等算子，在 \(F_{ij}\) 及不含这两个指标的基向量上为零；右侧恰好把前一类基向量变号，其余保持不变。Lie 理想被 \(\operatorname{ad}F_{ij}\) 及其多项式保持，故被每个 \(d_{ij}\) 共轭保持，从而被 \(D\) 共轭保持。

选取偶数个坐标变号，正是为了用这个内在的伴随算子多项式保证每个理想都被保持。一般 Lie 代数自同构虽然把理想送到理想，却未必保持一个指定理想。例如对非零 Lie 代数 \(\mathfrak s\)，\(\mathfrak s\oplus\mathfrak s\) 中交换两因子的自同构就把 \(\mathfrak s\oplus0\) 送到 \(0\oplus\mathfrak s\)。因此不能仅凭所有符号翻转都给出自同构，就把它们全部用于提取理想中的分量。

每个指标对 \(\{i,j\}\) 定义 \(D\) 的一维特征标 \(\chi_{ij}(d)=\varepsilon_i\varepsilon_j\)。对 \(n=3\) 或 \(n\ge5\)，不同二元集合给出不同特征标。事实上，两个这样的特征标相同，只可能是两集合相同，或它们的对称差为全部指标；后一个可能性恰好发生于 \(n=4\) 的互补二元集合。因为 \(|D|=2^{n-1}\) 在 \(\mathbb C\) 中非零，可以除以 \(|D|\) 取平均。相应特征空间的投影对 \(X=\sum_{i<j}x_{ij}F_{ij}\) 给出
\[
\frac1{|D|}\sum_{d=\operatorname{diag}(\varepsilon)\in D}
\varepsilon_i\varepsilon_j\,dXd^{-1}=x_{ij}F_{ij}.
\]
不相同的两个特征标相乘为非平凡特征标，其和为零：取它取值为 \(-1\) 的群元素平移求和即可。这正是有限交换群的一维特征标正交投影，在这里把一个理想中的矩阵分离为单个矩阵单位方向。例如在 \(\mathfrak{so}_3\) 中，\(D=\{1,d_{12},d_{13},d_{23}\}\cong C_2\times C_2\)。对 \(X=aF_{12}+bF_{13}+cF_{23}\)，共轭作用分别给出
\[
\begin{aligned}
\operatorname{Ad}_{d_{12}}X&=aF_{12}-bF_{13}-cF_{23},\\
\operatorname{Ad}_{d_{13}}X&=-aF_{12}+bF_{13}-cF_{23},\\
\operatorname{Ad}_{d_{23}}X&=-aF_{12}-bF_{13}+cF_{23}.
\end{aligned}
\]
因此相应 Fourier 投影为
\[
\begin{aligned}
\Pi_{12}(X)&=\frac14\bigl(X+\operatorname{Ad}_{d_{12}}X
-\operatorname{Ad}_{d_{13}}X-\operatorname{Ad}_{d_{23}}X\bigr)\\
&=aF_{12}.
\end{aligned}
\]
若 \(X\in I\) 且 \(a\ne0\)，各共轭作用与它们的线性组合都保持 \(I\)，故 \(F_{12}=a^{-1}\Pi_{12}(X)\in I\)。一般地，某个非零 \(x_{ij}\) 保证理想含有 \(F_{ij}\)。再用
\[
[F_{ij},F_{jk}]=F_{ik}\quad(i,j,k\;\text{互异})
\]
以及反对称记号 \(F_{ji}=-F_{ij}\)，先得到涉及 \(i\) 或 \(j\) 的全部基向量，再得到其余的 \(F_{k\ell}\)。所以非零理想等于 \(\mathfrak{so}_n\)。三类代数均非交换；其中心作为交换理想只能为零。
\end{proof}

上述正交 Lie 代数的证明在 \(n=4\) 时不能分离互补二元集合。例如 \(\chi_{12}=\chi_{34}\)，相应投影只能提取 \(\mathbb CF_{12}\oplus\mathbb CF_{34}\)，不能把这两个方向分开。特征标重合说明这个提取方法在四维时失效，本身还不能证明该代数不单；真正的理想分解需要如下括号计算。\(\mathfrak{so}_2\) 是一维交换代数，中心是自身；\(\mathfrak{so}_1=\mathfrak{sl}_1=0\)。

\begin{proposition}\label{b5:prop:so4-explicit-ideals}
在复 Lie 代数 \(\mathfrak{so}_4\) 中，记 \(F_{ij}=E_{ij}-E_{ji}\)，并令
\[
\begin{aligned}
a_1&=F_{12}+F_{34},&a_2&=F_{13}-F_{24},&a_3&=F_{14}+F_{23},\\
b_1&=F_{12}-F_{34},&b_2&=F_{13}+F_{24},&b_3&=F_{14}-F_{23}.
\end{aligned}
\]
则 \(\mathfrak a=\operatorname{span}_{\mathbb C}\{a_1,a_2,a_3\}\) 与 \(\mathfrak b=\operatorname{span}_{\mathbb C}\{b_1,b_2,b_3\}\) 是两个互相交换的三维单理想，且
\[
\mathfrak{so}_4=\mathfrak a\oplus\mathfrak b.
\]
因而 \(\mathfrak{so}_4\) 不单，但中心为零。
\end{proposition}
\begin{proof}
约定 \(F_{ji}=-F_{ij}\)、\(F_{ii}=0\)。由矩阵单位乘法得到
\[
[F_{ij},F_{k\ell}]
=\delta_{jk}F_{i\ell}-\delta_{ik}F_{j\ell}
-\delta_{j\ell}F_{ik}+\delta_{i\ell}F_{jk}.
\]
把六个向量的定义代入，得到
\[
\begin{aligned}
[a_1,a_2]&=-2a_3,&[a_2,a_3]&=-2a_1,&[a_3,a_1]&=-2a_2,\\
[b_1,b_2]&=2b_3,&[b_2,b_3]&=2b_1,&[b_3,b_1]&=2b_2,
\end{aligned}
\]
以及 \([a_i,b_j]=0\) 对全部 \(1\le i,j\le3\) 成立。例如
\[
\begin{aligned}
[a_1,a_2]&=-F_{23}-F_{14}-F_{14}-F_{23}=-2a_3,\\
[a_1,b_2]&=-F_{23}+F_{14}-F_{14}+F_{23}=0.
\end{aligned}
\]
六个向量由三对可逆的和差变换从 \(F_{ij}\) 得到，故构成一组基。上述括号使两个三维子空间各自闭合且彼此交换，所以它们是理想，并给出所述直和。基 \(a_1/2,a_2/2,a_3/2\) 与 \(-b_1/2,-b_2/2,-b_3/2\) 的循环括号都与 \(\mathfrak{so}_3\) 中 \(F_{12},F_{13},F_{23}\) 的相同，因此两个理想均同构于 \(\mathfrak{so}_3\)，由\cref{b5:prop:classical-simple}均为单 Lie 代数。它们的中心都为零，直和的中心也为零。
\end{proof}

四维情形还说明奇数次符号翻转为什么不能用于上述理想分量的投影。取 \(d=\operatorname{diag}(-1,1,1,1)\)，则 \(d\) 的共轭作用仍是 \(\mathfrak{so}_4\) 的自同构，但
\[
\operatorname{Ad}_d(a_i)=-b_i,\qquad
\operatorname{Ad}_d(b_i)=-a_i\quad(1\le i\le3).
\]
所以它交换 \(\mathfrak a\) 与 \(\mathfrak b\)，而不保持其中任意一个。扩大到所有符号翻转后，虽然可以区分 \(F_{12}\) 与 \(F_{34}\)，相应平均却未必留在原来的理想中。

\subsection{小秩同构与例外情形}
\label{b5:subsec:1304:03}

四维空间的二次外幂是六维的，两个二向量的楔积又落在一维最高外幂中，因而体积元可以把这个楔积解释为一个数。若原空间还带有辛型，其二次外幂中会多出一条由辛型决定的平凡直线；下面先把相应的五维补空间构造出来。

\begin{proposition}\label{b5:prop:sp4-primitive-wedge}
设 \(V\) 是四维复辛空间，取辛基 \(e_1,e_2,f_1,f_2\)，并令
\[
\operatorname{vol}=e_1\wedge e_2\wedge f_1\wedge f_2,\qquad
B(u,v)\operatorname{vol}=u\wedge v
\quad(u,v\in\bigwedge\nolimits^2V).
\]
则 \(B\) 是 \(\bigwedge^2V\) 上非退化的对称双线性型。令
\[
\Omega=e_1\wedge f_1+e_2\wedge f_2,\qquad W=\Omega^\perp.
\]
在 \(\mathfrak{sp}(V)\) 的外幂表示下，\(\mathbb C\Omega\) 是非退化的一维平凡子模，\(W\) 是五维非退化子模，并有子模直和
\[
\bigwedge\nolimits^2V=\mathbb C\Omega\oplus W.
\]
其中 \(W\) 的一组基为
\[
\begin{aligned}
w_1&=e_1\wedge e_2,&w_2&=e_1\wedge f_2,&w_3&=f_1\wedge e_2,\\
w_4&=f_1\wedge f_2,&w_5&=e_1\wedge f_1-e_2\wedge f_2.
\end{aligned}
\]
相对此辛基，令 \(H_i=E_{ii}-E_{2+i,2+i}\)。对 \(H=t_1H_1+t_2H_2\)，上述五个基向量的特征值依次为
\[
t_1+t_2,\quad t_1-t_2,\quad -t_1+t_2,\quad -t_1-t_2,\quad0.
\]
\end{proposition}
\begin{proof}
两个二向量满足 \(u\wedge v=(-1)^{2\cdot2}v\wedge u=v\wedge u\)，故 \(B\) 对称。对辛分块矩阵
\(X=\begin{psmallmatrix}A&B_0\\C&-A^{\mathsf T}\end{psmallmatrix}\)，其中 \(B_0,C\) 对称，直接作用于 \(\Omega\) 得
\[
\begin{aligned}
X\Omega={}&\sum_{i,j=1}^2
\bigl(A_{ji}e_j\wedge f_i-A_{ij}e_i\wedge f_j\bigr)\\
&+\sum_{i,j=1}^2 C_{ji}f_j\wedge f_i
+\sum_{i,j=1}^2 (B_0)_{ji}e_i\wedge e_j=0.
\end{aligned}
\]
第一行交换 \(i,j\) 后成对相消，第二行的两个和分别因对称性与楔积的反对称性相消。因此 \(\mathbb C\Omega\) 是平凡子模。

又有
\[
\Omega\wedge\Omega
=2e_1\wedge f_1\wedge e_2\wedge f_2=-2\operatorname{vol},
\qquad B(\Omega,\Omega)=-2\ne0.
\]
所以 \(B(\Omega,-)\) 是非零泛函，其核 \(W\) 为五维，并且这条直线与 \(W\) 的交为零。对前四个 \(w_i\)，与 \(\Omega\) 楔积时各项都有重复的基向量，故全为零。写 \(u=e_1\wedge f_1\)、\(v=e_2\wedge f_2\)，则
\[
w_5\wedge\Omega=(u-v)\wedge(u+v)
=u\wedge v-v\wedge u=0.
\]
因此五个 \(w_i\) 均在 \(W\) 中。它们线性无关，故为 \(W\) 的一组基。直接按体积元的次序计算配对，得到
\[
B(w_1,w_4)=1,\qquad B(w_2,w_3)=-1,\qquad B(w_5,w_5)=2,
\]
其余基向量配对除对称项外均为零。因此 \(B|_W\) 的 Gram 矩阵在基 \(w_1,w_4,w_2,w_3,w_5\) 下由两个可逆的二阶配对块与非零的一阶块组成，故非退化。这也说明 \(B\) 在整个直和 \(\mathbb C\Omega\oplus W=\bigwedge^2V\) 上非退化。

每个 \(X\in\mathfrak{sp}(V)\) 都迹零。由\cref{b5:prop:lie-top-exterior-trace}，它在 \(\bigwedge^4V\) 上的作用为零，所以外幂作用的 Leibniz 公式给出
\[
B(Xu,v)+B(u,Xv)=0\quad(u,v\in\bigwedge\nolimits^2V).
\]
对 \(w\in W\)，于是 \(B(Xw,\Omega)=-B(w,X\Omega)=0\)，故 \(W\) 被保持，上面的向量空间直和确为子模直和。最后，\(He_i=t_ie_i\)、\(Hf_i=-t_if_i\)，将两因子的特征值相加就得到所列五个值。
\end{proof}

\begin{proposition}[小维数同构]\label{b5:prop:classical-low-rank-isomorphisms}
复数域上有 Lie 代数同构
\[
\mathfrak{sp}_2\cong\mathfrak{sl}_2\cong\mathfrak{so}_3,\qquad
\mathfrak{so}_4\cong\mathfrak{sl}_2\oplus\mathfrak{sl}_2,
\qquad
\mathfrak{sp}_4\cong\mathfrak{so}_5,\qquad
\mathfrak{sl}_4\cong\mathfrak{so}_6.
\]
\end{proposition}
\begin{proof}
二维辛矩阵的块条件恰好是迹零，故 \(\mathfrak{sp}_2=\mathfrak{sl}_2\) 在标准基下直接相等。\(\mathfrak{sl}_2\) 的伴随表示保持非退化对称 Killing 型；其核是零中心。因此它给出单射
\(\mathfrak{sl}_2\rightarrow\mathfrak{so}(\mathfrak{sl}_2,\kappa)\)，两边维数均为 \(3\)，便得到与 \(\mathfrak{so}_3\) 的同构。

设 \(U,W\) 都是带非退化交替型的二维复空间。两个交替型的张量积是 \(U\otimes W\) 上的非退化对称型，且
\[
(A,B)\longmapsto A\otimes I+I\otimes B
\]
把 \(\mathfrak{sl}(U)\oplus\mathfrak{sl}(W)\) 映入相应的四维正交代数。这里通过自然同构 \(\operatorname{End}(U\otimes W)\cong\operatorname{End}(U)\otimes\operatorname{End}(W)\)，定义对第二因子的部分迹为
\[
1\otimes\operatorname{tr}:\operatorname{End}(U)\otimes\operatorname{End}(W)
\longrightarrow\operatorname{End}(U),\qquad
C\otimes D\longmapsto(\operatorname{tr}D)C.
\]
若像为零，应用这个线性映射得到 \(2A+(\operatorname{tr}B)I_U=2A=0\)。同理，对第一因子取部分迹得到 \((\operatorname{tr}A)I_W+2B=2B=0\)。因此映射单射。两边维数为 \(6\)，故得到 \(\mathfrak{so}_4\) 的分解。

在四维复空间 \(V\) 上取基 \(e_1,e_2,e_3,e_4\)，并令 \(\operatorname{vol}=e_1\wedge e_2\wedge e_3\wedge e_4\)，以
\[
B(u,v)\operatorname{vol}=u\wedge v,\qquad u,v\in\bigwedge\nolimits^2V
\]
定义六维空间上的型。两个二向量交换时符号为 \((-1)^{2\cdot2}=1\)，故这个型对称。记 \(e_{ij}=e_i\wedge e_j\)，则
\[
e_{12}\wedge e_{34}=\operatorname{vol},\qquad
e_{13}\wedge e_{24}=-\operatorname{vol},\qquad
e_{14}\wedge e_{23}=\operatorname{vol}.
\]
除了这三对及其交换次序，两个基二向量的楔积都有重复指标而为零。因此在基
\[
e_{12},e_{34},e_{13},-e_{24},e_{14},e_{23}
\]
下，Gram 矩阵由三个 \(\begin{psmallmatrix}0&1\\1&0\end{psmallmatrix}\) 对角块组成，故 \(B\) 非退化。对 \(X\in\mathfrak{sl}(V)\)，\cref{b5:prop:lie-top-exterior-trace}给出 \(X\operatorname{vol}=0\)，于是
\[
\bigl(B(Xu,v)+B(u,Xv)\bigr)\operatorname{vol}
=X(u\wedge v)=0.
\]
所以 \(\mathfrak{sl}(V)\) 保持 \(B\)。对于迹零对角矩阵 \(H=\operatorname{diag}(t_1,t_2,t_3,t_4)\)，有 \(He_{ij}=(t_i+t_j)e_{ij}\)；互补的一对指标满足 \(t_i+t_j=-(t_k+t_\ell)\)，所以三对非零配对方向具有相反的权。这是一般不变性在对角作用中的具体表现。这个表示非零，例如
\(\operatorname{diag}(1,0,-1,0)\) 在 \(e_1\wedge e_2\) 上非零。由\cref{b5:prop:classical-simple}，单 Lie 代数 \(\mathfrak{sl}_4\) 的非零表示核为零，故得到单射 \(\mathfrak{sl}_4\rightarrow\mathfrak{so}_6\)。两边均为 \(15\) 维，因而同构。

最后令四维 \(V\) 带辛型。\cref{b5:prop:sp4-primitive-wedge}说明，同一个六维外幂空间作为 \(\mathfrak{sp}_4\)-表示分解为一维平凡子模 \(\mathbb C\Omega\) 与五维非退化正交子模 \(W\)。因此得到 Lie 代数同态
\[
\mathfrak{sp}_4\longrightarrow\mathfrak{so}(W,B|_W).
\]
这个表示非零：上述 \(w_1=e_1\wedge e_2\in W\)，而 \(H_1w_1=w_1\)。辛代数的单性使其核为零；源与 \(\mathfrak{so}(W,B|_W)\) 均为 \(10\) 维，所以同态为同构，得到 \(\mathfrak{sp}_4\cong\mathfrak{so}_5\)。这里的五维空间是辛型决定的平凡直线的正交补，整个六维外幂表示则是 \(\mathbf1\oplus W\)，其中 \(\mathbf1\) 表示一维平凡表示。
\end{proof}

这些同构是 Lie 代数层面的同构。相应矩阵群还可能相差有限中心或覆盖，不能把本命题直接改写成同名群的同构。后续讨论自旋群和代数群时，会进一步区分这两种结论。


\begin{chaptersummary}
	迹型的不变性来自矩阵迹的循环性。Cartan 可解性判据进一步说明，当导出代数与整个代数在迹型下正交时，导出代数的每个算子都幂零。证明中引入共轭特征值构造，排除了复数平方相消的可能，再由 Engel 定理完成推导。

	\ProseParagraphBreak

	Killing 型非退化与半单性等价。理想的正交补因此成为可计算的理想直和补，所有单分量以及所有理想都被确定。这个分解针对伴随表示，不能在未证明 Weyl 定理之前推广成任意表示的完全可约性。

	\ProseParagraphBreak

	经典矩阵族的单性已由矩阵作用直接证明：特殊线性与辛情形用对角伴随作用提取权分量，正交情形用偶数次符号翻转提取反对称矩阵单位。\(\mathfrak{so}_4\) 的例外对应两个单理想，其他小维数同构则由伴随表示、张量积和外幂构造实现。根系理论将把这些矩阵现象组织成统一的分类语言。
\end{chaptersummary}

\BookExercises

\begin{exercise}\label{b5:ex:13-killing-basis}
	在 \(\mathfrak{sl}_2(\mathbb C)\) 中，求 Killing 型相对于 \(h,e+f,e-f\) 的矩阵。由此判断 \(\mathbb C(e+f)\) 和 \(\mathbb Ce\) 上的限制是否退化。
\end{exercise}
\begin{solution}
	由 \(\kappa(h,h)=8\)、\(\kappa(e,f)=4\)，新基下的矩阵为 \(\operatorname{diag}(8,8,-8)\)，所有不同基向量的配对为零。因此 \(\mathbb C(e+f)\) 上的限制非退化，而 \(\kappa(e,e)=0\) 使 \(\mathbb Ce\) 上的限制恒零。这不与全空间非退化矛盾：非退化双线性型可以含有各向同性向量。
\end{solution}

\begin{exercise}\label{b5:ex:13-gln-killing-radical}
	对 \(n\ge2\)，计算 \(\mathfrak{gl}_n(\mathbb C)\) 的 Killing 型的根，并将其与该代数的中心比较。
\end{exercise}
\begin{solution}
	正文给出 \(\kappa(X,Y)=2n\operatorname{tr}(XY)-2\operatorname{tr}X\,\operatorname{tr}Y\)。 标量矩阵显然与全部矩阵正交。写 \(X=X_0+\operatorname{tr}(X)I/n\)，其中 \(\operatorname{tr}X_0=0\)。若 \(X\) 在型的根中，则 \(\operatorname{tr}(X_0Y)=0\) 对所有迹零 \(Y\) 成立。依次取 \(Y=E_{ji}\) 消去 \(X_0\) 的非对角元，再取 \(Y=E_{ii}-E_{jj}\) 说明对角元相等；迹零使其全零。因此根恰为 \(\mathbb CI\)，等于中心。
\end{solution}

\begin{exercise}\label{b5:ex:13-heisenberg-killing}
	计算三维复 Heisenberg Lie 代数 \([x,y]=z\) 的全部伴随矩阵及 Killing 型。再证明任意有限维幂零 Lie 代数的 Killing 型恒零。
\end{exercise}
\begin{solution}
	在基 \(x,y,z\) 下，\(\operatorname{ad}x=E_{32}\)、 \(\operatorname{ad}y=-E_{31}\)、\(\operatorname{ad}z=0\)；任意两个的乘积为零，所以 Killing 型恒零。一般地，若 \(\mathfrak g\) 幂零，则对任意 \(a\in\mathfrak g\) 和整数 \(m\ge0\) 有
	\[
		(\operatorname{ad}a)^m(\mathfrak g)\subseteq\gamma_{m+1}(\mathfrak g).
	\]
	下降中心列最终为零，故每个 \(\operatorname{ad}a\) 都幂零。Engel 定理给出共同的严格上三角基；两个这种矩阵的乘积仍严格上三角，迹为零。因此全部 Killing 配对都为零。
\end{solution}

\begin{exercise}\label{b5:ex:13-dual-trace}
	设 \(\rho\) 是有限维复 Lie 代数 \(\mathfrak g\) 的有限维表示，\(\rho^\vee\) 是对偶表示。证明两者给出的迹型相同。
\end{exercise}
\begin{solution}
	在对偶基下，\(\rho^\vee(x)=-\rho(x)^{\mathsf T}\)。因此
	\[
		\operatorname{tr}\bigl(\rho^\vee(x)\rho^\vee(y)\bigr) =\operatorname{tr}\bigl(\rho(x)^{\mathsf T}\rho(y)^{\mathsf T}\bigr) =\operatorname{tr}\bigl(\rho(y)\rho(x)\bigr) =\operatorname{tr}\bigl(\rho(x)\rho(y)\bigr).
	\]
	负号成对抵消，而转置把乘法次序反转，最后用迹的循环性。
\end{solution}

\begin{exercise}\label{b5:ex:13-sl2-invariant-direct}
	不使用不变型唯一性命题，直接由标准括号计算 \(\mathfrak{sl}_2(\mathbb C)\) 上全部对称不变双线性型。
\end{exercise}
\begin{solution}
	对 \([h,e]=2e\)、\([h,f]=-2f\) 使用不变性，得到 \(B(e,e)=B(f,f)=B(h,e)=B(h,f)=0\)。例如 \(B\bigl([h,e],e\bigr)=B\bigl(h,[e,e]\bigr)=0\)，所以 \(2B(e,e)=0\)； \(B\bigl([h,h],e\bigr)=B\bigl(h,[h,e]\bigr)\) 给出 \(B(h,e)=0\)。 再由
	\[
		B(h,h)=B\bigl([e,f],h\bigr)=B\bigl(e,[f,h]\bigr)=2B(e,f)
	\]
	可知只剩一个自由参数 \(c=B(e,f)\)。矩阵为
	\(\begin{psmallmatrix}2c&0&0\\0&0&c\\0&c&0\end{psmallmatrix}\)，这里基的次序为 \(h,e,f\)。这些矩阵直接代入三条基括号均满足不变性，故恰为全部解，即 \((c/4)\kappa\)。
\end{solution}

\begin{exercise}\label{b5:ex:13-subalgebra-killing}
	取 \(\mathfrak h=\mathbb Ch\subseteq\mathfrak{sl}_2(\mathbb C)\)。比较 \(\kappa_{\mathfrak h}\) 与 \(\kappa_{\mathfrak{sl}_2}|_{\mathfrak h}\)，指出理想限制引理中哪个条件不能删去。
\end{exercise}
\begin{solution}
	\(\mathfrak h\) 交换，所以自身伴随表示为零，\(\kappa_{\mathfrak h}=0\)。但 \(\kappa_{\mathfrak{sl}_2}(h,h)=8\)，环境 Killing 型的限制非零。这里 \(\mathfrak h\) 不是理想，因为 \([e,h]=-2e\notin\mathfrak h\)。因此 \cref{b5:lem:killing-restriction-ideal}的“理想”不能换成“子代数”。
\end{solution}

\begin{exercise}\label{b5:ex:13-polarized-trace}
	设 \(L\subseteq\mathfrak{gl}(V)\) 是有限维复矩阵 Lie 代数，并且 \(\operatorname{tr}(x^2)=0\) 对全部 \(x\in L\) 成立。证明 \(L\) 可解。再构造非幂零的三阶矩阵 \(x\)，使 \(\operatorname{tr}x=\operatorname{tr}(x^2)=0\)，并计算 \(\operatorname{tr}(x^3)\)。
\end{exercise}
\begin{solution}
	将 \(x+y\) 代入并相减，得 \(0=\operatorname{tr}(xy+yx)=2\operatorname{tr}(xy)\)。 因此迹型在全部 \(L\times L\) 上为零，特别满足 Cartan 可解性判据，故 \(L\) 可解。取满足 \(\omega^2+\omega+1=0\) 的复数 \(\omega\)，令 \(x=\operatorname{diag}(1,\omega,\omega^2)\)。因 \(\omega^3=1\)，有
	\[
		\operatorname{tr}x=1+\omega+\omega^2=0,\qquad \operatorname{tr}(x^2)=1+\omega^2+\omega^4=0,\qquad \operatorname{tr}(x^3)=3.
	\]
	而 \(x^3=I\)，所以 \(x\) 不幂零。前两项迹为零仍不足以排除非零特征值；Cartan 判据使用的是整个代数中的配对条件。
\end{solution}

\begin{exercise}\label{b5:ex:13-zero-killing-not-nilpotent}
	给定 \(a,b\in\mathbb C\)，在复空间 \(\mathbb Ch\oplus\mathbb Cx\oplus\mathbb Cy\) 上规定 \([h,x]=ax\)、\([h,y]=by\)、\([x,y]=0\)，得到 Lie 代数 \(\mathfrak g_{a,b}\)。求它的可解根、幂零根及 Killing 型的根，并确定 Killing 型恒零和代数幂零各自的参数条件。
\end{exercise}
\begin{solution}
	唯一需要核对的三个不同基向量的 Jacobi 恒等式中，各项都为零，故这些括号确实定义 Lie 代数。在基 \(h,x,y\) 下，伴随矩阵为
	\[
		\operatorname{ad}h=\operatorname{diag}(0,a,b),\quad \operatorname{ad}x=-aE_{21},\quad \operatorname{ad}y=-bE_{31}.
	\]
	所以 Killing 矩阵为 \(\operatorname{diag}(a^2+b^2,0,0)\)。记 \(V=\mathbb Cx\oplus\mathbb Cy\)，则
	\[
		\operatorname{rad}(\kappa)=
		\begin{cases}
			V,&a^2+b^2\ne0,\\
			\mathfrak g_{a,b},&a^2+b^2=0.
		\end{cases}
	\]
	由于导出代数包含于交换理想 \(V\)，代数始终可解，可解根就是整个 \(\mathfrak g_{a,b}\)。

	若 \(a=b=0\)，代数交换，幂零根也是整个代数。否则 \(V\) 是幂零理想，且任何幂零理想都包含于 \(V\)：假设理想 \(J\) 含有 \(u=ch+v\)，其中 \(c\ne0\)、\(v\in V\)。若 \(a\ne0\)，则 \([u,x]=cax\) 迫使 \(x\in J\)，而 \((\operatorname{ad}_J u)^m x=(ca)^m x\ne0\)，与 \(J\) 幂零矛盾；若 \(b\ne0\)，对 \(y\) 作同样计算。因此非零参数情形的幂零根恰为 \(V\)，代数本身不幂零。综上，Killing 型恒零当且仅当 \(a^2+b^2=0\)，而代数幂零当且仅当 \(a=b=0\)。
\end{solution}

\begin{exercise}\label{b5:ex:13-two-dimensional-killing}
	证明每个二维复 Lie 代数的 Killing 型都退化。
\end{exercise}
\begin{solution}
	交换情形的 Killing 型为零。非交换时，取基 \(u,v\) 并令 \(e=[u,v]=au+bv\ne0\)。有 \([u,e]=be\)、\([v,e]=-ae\)。可取 \(h=\alpha u+\beta v\) 使 \(\alpha b-\beta a=1\)，于是 \([h,e]=e\)，且 \(h,e\) 线性无关。故代数同构于正文的二维仿射例子，Killing 矩阵为 \(\operatorname{diag}(1,0)\)，仍退化。
\end{solution}

\begin{exercise}\label{b5:ex:13-trace-nonfaithful}
	设 \(\mathfrak g=\mathfrak{sl}_2(\mathbb C)\oplus\mathfrak{sl}_2(\mathbb C)\)。在四维空间 \(V=\mathbb C^2\oplus\mathbb C^2\) 与 \(W=\mathbb C^2\otimes\mathbb C^2\) 上分别取表示
	\[
		\rho_{\oplus}(X,Y)=\operatorname{diag}(X,Y),\qquad \rho_{\otimes}(X,Y)=X\otimes I+I\otimes Y.
	\]
	证明两者忠实，计算各自的迹型，并判断它们是否同构为 \(\mathfrak g\)-模。
\end{exercise}
\begin{solution}
	第一种表示的核显然为零。第二种若满足 \(X\otimes I+I\otimes Y=0\)，对第二因子取部分迹得到 \(2X+(\operatorname{tr}Y)I=2X=0\)，对第一因子取部分迹得到 \(2Y=0\)，所以也忠实。 分块乘法给出
	\[
		\beta_{\oplus}\bigl((X,Y),(X',Y')\bigr) =\operatorname{tr}(XX')+\operatorname{tr}(YY').
	\]
	张量表示的乘积除 \(XX'\otimes I\) 和 \(I\otimes YY'\) 外，还有 \(X\otimes Y'\) 和 \(X'\otimes Y\)。利用 \(\operatorname{tr}(A\otimes B)=\operatorname{tr}A\,\operatorname{tr}B\)，交叉项因 \(X,Y,X',Y'\) 都迹零而消失，得到
	\[
		\beta_{\otimes}\bigl((X,Y),(X',Y')\bigr) =2\operatorname{tr}(XX')+2\operatorname{tr}(YY').
	\]
	由于 \(\kappa_{\mathfrak{sl}_2}(X,X')=4\operatorname{tr}(XX')\)，两种迹型分别是直和 Killing 型的 \(1/4\) 倍与 \(1/2\) 倍，均非退化。同构表示的作用矩阵同时共轭，因而迹型必相同；这里两型不同，例如在 \(((h,0),(h,0))\) 上分别取值 \(2\) 与 \(4\)，故两个四维模并不同构。
\end{solution}

\begin{exercise}\label{b5:ex:13-count-ideals}
	设 \(\mathfrak g=\bigoplus_{i=1}^r\mathfrak g_i\) 是有限维复半单 Lie 代数的单理想分解。证明它恰有 \(2^r\) 个理想，并确定每个商的单分量。
\end{exercise}
\begin{solution}
	\cref{b5:thm:semisimple-ideal-decomposition}把每个理想唯一写成 \(I_S=\bigoplus_{i\in S}\mathfrak g_i\)，其中 \(S\subseteq\{1,\ldots,r\}\)。不同子集给出不同子空间，故理想数为 \(2^r\)。 商映射在其余分量的直和上为同构，所以 \(\mathfrak g/I_S\cong\bigoplus_{i\notin S}\mathfrak g_i\)。
\end{solution}

\begin{exercise}\label{b5:ex:13-automorphisms-components}
	设 \(\mathfrak s\) 是非零有限维复单 Lie 代数。证明
	\[
		\operatorname{Aut}(\mathfrak s^{\oplus r}) \cong\operatorname{Aut}(\mathfrak s)^r\rtimes S_r,
	\]
	其中 \(S_r\) 置换各因子。
\end{exercise}
\begin{solution}
	自同构必置换极小非零理想，而这些理想恰为 \(r\) 个坐标因子，因此给出到 \(S_r\) 的同态。每个坐标置换确实是自同构，提供该同态的分裂。核中的自同构逐个保持每个坐标理想，其限制是 \(r\) 个独立的 \(\mathfrak s\) 自同构；反之这些自同构逐坐标作用便给出核的所有元素。共轭作用正是置换这 \(r\) 个自同构，因此为所述半直积。
\end{solution}

\begin{exercise}\label{b5:ex:13-all-invariant-forms-semisimple}
	设 \(\mathfrak s\) 是非零有限维复单 Lie 代数，\(r\ge2\)。在 \(\mathfrak g=\mathfrak s^{\oplus r}\) 的对称不变双线性型中，求出同时被全部 Lie 代数自同构保持的型，并判断何时非退化。
\end{exercise}
\begin{solution}
	由\cref{b5:prop:semisimple-invariant-forms}，可写
	\[
		B=\bigoplus_{i=1}^r c_i\kappa_{\mathfrak s}.
	\]
	交换任意两个坐标因子都是自同构。若 \(B\) 被全部自同构保持，对第一、第二坐标中相同的非零 Killing 配对比较，便得到 \(c_1=c_2\)；对任意一对坐标同理，故全部系数相同。于是只能有 \(B=c\kappa_{\mathfrak g}\)。 反过来，若 \(\varphi\) 是任意自同构，则 \(\operatorname{ad}(\varphi x)=\varphi\operatorname{ad}(x)\varphi^{-1}\)，所以
	\[
		\kappa_{\mathfrak g}(\varphi x,\varphi y) =\operatorname{tr}\bigl(\varphi\operatorname{ad}x\operatorname{ad}y\varphi^{-1}\bigr) =\kappa_{\mathfrak g}(x,y).
	\]
	因此 Killing 型及其每个倍数都被全部自同构保持。因 \(\mathfrak g\) 半单且非零，这些型非退化当且仅当 \(c\ne0\)。仅要求 Lie 不变性时可独立选择各系数，增加自同构不变性才把它们约束为同一个数。
\end{solution}

\begin{exercise}\label{b5:ex:13-diagonal-perpendicular}
	在 \(\mathfrak s\oplus\mathfrak s\) 中，\(\mathfrak s\) 是非零复单 Lie 代数。计算对角子代数的中心化子与 Killing 正交补，判断这个正交补是否为子代数。
\end{exercise}
\begin{solution}
	若 \((a,b)\) 中心化全部 \((x,x)\)，则 \([a,x]=[b,x]=0\) 对所有 \(x\) 成立，故 \(a=b=0\)。中心化子为零。 另一方面，直和 Killing 型给出
	\[
		\kappa\bigl((a,b),(x,x)\bigr)=\kappa_{\mathfrak s}(a+b,x),
	\]
	所以正交补为 \(\bigl\{\,(a,-a)\mid a\in\mathfrak s\,\bigr\}\)。 取 \([a,b]\ne0\)，则两反对角元素的括号为 \(\bigl([a,b],[a,b]\bigr)\)，不在反对角空间中。因此这个正交补不是子代数，更不是理想。
\end{solution}

\begin{exercise}\label{b5:ex:13-sln-center-elementary}
	不使用单性，直接证明 \(n\ge2\) 时 \(\mathfrak{sl}_n(\mathbb C)\) 的中心为零。指出在特征 \(p\mid n\) 时哪一步会失效。
\end{exercise}
\begin{solution}
	若 \(X\) 与全部迹零对角矩阵交换，则 \(X\) 的非对角元为零，因为对每对 \(i\ne j\) 可选一组迹零对角元使第 \(i,j\) 项不同。再与所有 \(E_{ij}\) 交换，便迫使所有对角元相等，即 \(X=cI\)。在复数域上 \(\operatorname{tr}X=nc=0\) 给出 \(c=0\)。 若 \(\operatorname{char}k=p\mid n\)，由 \(nc=0\) 不能再推出 \(c=0\)；事实上 \(I\in\mathfrak{sl}_n(k)\) 是非零中心元素。在小特征下前面的对角分离也可能需要调整，但这个标量反例本身已经足够。
\end{solution}

\begin{exercise}\label{b5:ex:13-sl3-centralizer}
	求 \(\mathfrak{sl}_3(\mathbb C)\) 中 \(E_{12}\) 的中心化子及其维数，并说明它为什么不是整个 Lie 代数的理想。
\end{exercise}
\begin{solution}
	令 \(X=(x_{ij})\)。条件 \(XE_{12}=E_{12}X\) 给出 \(x_{21}=x_{23}=x_{31}=0\)、\(x_{11}=x_{22}\)，结合迹零，
	\[
		X=\begin{pmatrix}a&b&c\\0&a&0\\0&d&-2a\end{pmatrix}.
	\]
	故中心化子四维。它含有 \(E_{12}\)，但 \([E_{21},E_{12}]=E_{22}-E_{11}\) 不中心化 \(E_{12}\)，所以不是理想。这里的中心化子针对一个元素，并不适用半单理想的正交补结论。
\end{solution}

\begin{exercise}\label{b5:ex:13-sp4-basis}
	对 \(\mathfrak{sp}_4(\mathbb C)\)，列出正文基 \(H_i,A_{ij},B_{ij},C_{ij}\) 的全部十个元素及其对 \(H=t_1H_1+t_2H_2\) 的权。
\end{exercise}
\begin{solution}
	零权基为 \(H_1,H_2\)。其余为
	\[
		\begin{array}{c|cccccccc}
			\text{向量}&A_{12}&A_{21}&B_{11}&B_{22}&C_{11}&C_{22}&B_{12}&C_{12}\\
			\hline
			\text{权}&t_1-t_2&t_2-t_1&2t_1&2t_2&-2t_1&-2t_2&t_1+t_2&-t_1-t_2
		\end{array}
	\]
	其中 \(\bar1=3,\bar2=4\)，例如 \(A_{12}=E_{12}-E_{43}\)、\(B_{12}=E_{14}+E_{23}\)、 \(C_{12}=E_{32}+E_{41}\)。其余单指标元素由正文公式直接写出。这十个元素分别占据 \(A\) 块的四个自由位置及两个对称块各三个自由位置，因而线性无关且张成全部代数。
\end{solution}

\begin{exercise}\label{b5:ex:13-so4-explicit}
	取 \(d=\operatorname{diag}(-1,1,1,1)\)，令 \(\sigma(X)=dXd^{-1}\) 作用于 \(\mathfrak{so}_4(\mathbb C)\)。求 \(\sigma\) 的两个特征空间，判断不动子代数是否为理想，并用\cref{b5:prop:so4-explicit-ideals}的两个单理想解释它在直和分解中的位置。
\end{exercise}
\begin{solution}
	\(d\) 的共轭作用使 \(F_{1j}\) 变号，其余 \(F_{ij}\) 不变。因此两个特征空间为
	\[
		\mathfrak g_+=\operatorname{span}\{F_{23},F_{24},F_{34}\},\qquad \mathfrak g_-=\operatorname{span}\{F_{12},F_{13},F_{14}\}.
	\]
	\(\mathfrak g_+\) 是作用在最后三个坐标上的 \(\mathfrak{so}_3\)，所以为单子代数，但不是理想：例如 \([F_{12},F_{23}]=F_{13}\notin\mathfrak g_+\)。自同构的特征空间满足 \([\mathfrak g_+,\mathfrak g_+]\subseteq\mathfrak g_+\)、 \([\mathfrak g_+,\mathfrak g_-]\subseteq\mathfrak g_-\)、 \([\mathfrak g_-,\mathfrak g_-]\subseteq\mathfrak g_+\)，并不因此都是理想。

	在正文的基下，直接代入得 \(\sigma(a_i)=-b_i\)、\(\sigma(b_i)=-a_i\)。记两个单理想为 \(I_+=\operatorname{span}\{a_i\}\) 和 \(I_-=\operatorname{span}\{b_i\}\)，则 \(\varphi=\sigma|_{I_+}:I_+\to I_-\) 是 Lie 代数同构。不动子代数为
	\[
		\mathfrak g_+=\{\,u+\varphi(u)\mid u\in I_+\,\}.
	\]
	所以它是在两个单理想之间取同构图得到的对角型子代数，并非单理想分解中的一个因子。
\end{solution}

\begin{exercise}\label{b5:ex:13-even-sign-exception}
	分别在 \(\mathfrak{so}_4(\mathbb C)\) 与 \(\mathfrak{so}_5(\mathbb C)\) 中取 \(X=F_{12}+F_{34}\)。求包含 \(X\) 的最小 Lie 理想及其维数。在五维情形中，显式使用第 \(1,5\) 坐标同时变号的共轭作用，从 \(X\) 提取 \(F_{12}\)。
\end{exercise}
\begin{solution}
	四维时 \(X=a_1\) 属于\cref{b5:prop:so4-explicit-ideals}中的单理想 \(I_+=\operatorname{span}\{a_1,a_2,a_3\}\)。由 \([a_1,a_2]=-2a_3\) 和 \([a_3,a_1]=-2a_2\)，任何包含 \(X\) 的理想都包含这三项，故最小理想恰为 \(I_+\)，维数为 \(3\)。

	五维时令 \(d_{15}=\operatorname{diag}(-1,1,1,1,-1)\)。偶数次变号的共轭作用是 \(\operatorname{ad}F_{15}\) 的多项式，故保持每个理想。若 \(I\) 是包含 \(X\) 的理想，则
	\[
		\operatorname{Ad}_{d_{15}}X=-F_{12}+F_{34},\qquad F_{12}=\frac{X-\operatorname{Ad}_{d_{15}}X}{2}\in I.
	\]
	对 \(j=3,4,5\)，用 \([F_{12},F_{2j}]=F_{1j}\) 得 \(F_{1j}\in I\)。再由 \([F_{1j},F_{1k}]=-F_{jk}\) 得到所有其余矩阵单位，所以 \(I=\mathfrak{so}_5\)，维数为 \(10\)。同一个表达式在四维只生成一个单因子，在五维则生成整个代数；多出的坐标使这次理想内的分量分离成为可能。
\end{solution}

\begin{exercise}\label{b5:ex:13-primitive-wedge-five}
	在\cref{b5:prop:sp4-primitive-wedge}的五维空间 \(W\) 及其基的记号下，计算 \(B_{12}=E_{14}+E_{23}\) 的作用矩阵和 Jordan 块。比较 \(B_{12}\) 在 \(V=\mathbb C^4\) 上与在 \(W\) 上的幂零指数，并解释为什么不同的幂零指数不与 Lie 表示的定义矛盾。
\end{exercise}
\begin{solution}
	依次把正文五个基向量记为 \(w_1,\ldots,w_5\)。令 \(B=B_{12}\)，则在 \(V\) 上有 \(Bf_1=e_2\)、\(Bf_2=e_1\)、\(Be_1=Be_2=0\)，所以 \(B\ne0\) 而 \(B^2=0\)。在外幂上使用 \(\rho(B)(u\wedge v)=Bu\wedge v+u\wedge Bv\)，得到
	\[
		\rho(B)w_1=\rho(B)w_2=\rho(B)w_3=0,\qquad \rho(B)w_4=-w_5,\qquad \rho(B)w_5=2w_1.
	\]
	例如 \(\rho(B)(f_1\wedge f_2)=e_2\wedge f_2+f_1\wedge e_1=-w_5\)。因此在这组基下，作用矩阵为
	\[
		\begin{pmatrix}
			0&0&0&0&2\\
			0&0&0&0&0\\
			0&0&0&0&0\\
			0&0&0&0&0\\
			0&0&0&-1&0
		\end{pmatrix}.
	\]
	链 \(w_4\mapsto-w_5\mapsto-2w_1\mapsto0\) 给出一个三阶幂零 Jordan 块，而 \(w_2,w_3\) 分别给出一阶零块。故 Jordan 形为 \(J_3(0)\oplus J_1(0)\oplus J_1(0)\)，幂零指数为 \(3\)，大于 \(V\) 上的指数 \(2\)。Lie 表示要求保持括号，并不要求保持矩阵乘积；因此不能由 \(B^2=0\) 推出 \(\rho(B)^2=0\)。
\end{solution}

\begin{exercise}\label{b5:ex:13-exterior-trace-index}
	设 \(V=\mathbb C^4\)，\(\rho\) 是 \(\mathfrak{sl}(V)\) 在 \(\bigwedge^2V\) 上的表示。证明
	\[
		\operatorname{tr}_{\wedge^2V}\bigl(\rho(X)\rho(Y)\bigr) =2\operatorname{tr}_V(XY)\qquad(X,Y\in\mathfrak{sl}_4).
	\]
\end{exercise}
\begin{solution}
	两侧都是 \(\mathfrak{sl}_4\) 上的对称不变型，单性与 \cref{b5:cor:simple-killing-forms}说明它们相差一个标量。因此只须用一对非零配对确定标量。取 \(X=Y=\operatorname{diag}(1,-1,0,0)\)。在六个楔积基上，\(\rho(X)\) 的特征值为 \(0,1,1,-1,-1,0\)，平方和为 \(4\)，而 \(\operatorname{tr}_V(X^2)=2\)。比例为 \(2\)，得到结论。这里比较的是自然表示和外幂表示的迹型，不能省去这个比例而说两个迹型完全相同。
\end{solution}
